\chapter{Hodge 理论与 Yukawa 耦合}\label{chap:hodge-yukawa}
\SourcePage{96}

本章回顾并发展 Hodge 理论中的一些概念，将其用于计算 Calabi--Yau 流形 $V$ 的 Yukawa 耦合。为此，必须先理解 $V$ 上的周期积分所满足的 Picard--Fuchs 微分方程。我们将说明求出这些方程的两种方法。

\hyperref[sec:hodge-theory]{第 5.1 节}回顾 Hodge 结构变分、Picard--Fuchs 方程以及边界行为；\hyperref[sec:maximal-unipotent]{第 5.2 节}定义\emph{最大幺幂边界点}，它们在确定镜像映射时起重要作用。\hyperref[sec:griffiths-dwork]{第 5.3 节}解释 Griffiths 与 Dwork 发展的方法，用射影空间上有理微分形式的留数表示 $V$ 上的上同调类。\hyperref[sec:hodge-examples]{第 5.4 节}给出两个例题。\hyperref[sec:hypergeometric]{第 5.5 节}研究计算 Picard--Fuchs 方程的另一种方法，涉及 Gelfand、Kapranov 与 Zelevinsky 的\emph{广义超几何方程}。然后在\hyperref[sec:yukawa-couplings]{第~\ref*{sec:yukawa-couplings} 节}将这些内容用于计算 Yukawa 耦合。

本章有两个出人意料的方面：\hyperref[subsec:nef-gkz]{第~\ref*{subsec:nef-gkz} 小节}的 GKZ 方程组以上同调为值的解 $\widetilde I$，将在\hyperref[chap:mirror-theorem]{第~\ref*{chap:mirror-theorem} 章}的镜像定理讨论中出现；\hyperref[subsec:normalized-yukawa-connection]{第~\ref*{subsec:normalized-yukawa-connection} 小节}归一化 Yukawa 耦合与 Gauss--Manin 联络之间的关系，将在\hyperref[chap:quantum-cohomology]{第~\ref*{chap:quantum-cohomology} 章}讨论的 A 模型 Hodge 结构变分中发挥突出作用。

\section{Hodge 理论}\label{sec:hodge-theory}

本节回顾镜像对称中使用的 Hodge 理论。我们遵循 \SourceCite{Morrison2} 的思路，同时推广到参数空间维数更高的族。关于 Hodge 结构变分的一般参考文献有 \SourceCite{Griffiths1,Griffiths2,Schmid,CaK}；Calabi--Yau 变分的具体内容可参见 \SourceCite{BG}。

\subsection{Hodge 结构及其变分}\label{subsec:hodge-variations}

若 $X$ 是 $n$ 维光滑复射影簇，则 $H^k(X,\mathbb C)$ 有 Hodge 分解
\[
H^k(X,\mathbb C)\simeq\bigoplus_{p+q=k}H^{p,q}(X),
\]
其中相对于 $H^k(X,\mathbb R)$ 所确定的实结构，有 $\overline{H^{p,q}(X)}=H^{q,p}(X)$。再加入来自 $H^k(X,\mathbb Z)$ 的整数格，就得到一个\emph{权为 $k$ 的 Hodge 结构}。

我们常用 \emph{Hodge 滤过} $F^\bullet(X)=\{F^p(X)\}_{p=0}^n$ 来表述 $H^k(X,\mathbb C)$ 上的 Hodge 结构，其中
\[
F^p(X)=\bigoplus_{a\geq p}H^{a,k-a}(X).
\]
\SourcePage{97}
注意，$F^p(X)\oplus\overline{F^{k-p+1}(X)}=H^k(X,\mathbb C)$，且
\[
H^{p,q}(X)=F^p(X)\cap\overline{F^q(X)}.
\]
还有 $H^{p,q}(X)\simeq F^p(X)/F^{p+1}(X)$，其中 $p+q=k$。因此，$H^k(X,\mathbb C)$ 上的 Hodge 结构由整数结构与 Hodge 滤过唯一确定。

最后，$X$ 上的 Kähler 类 $\omega$ 在 $H^k(X,\mathbb C)$ 中确定\emph{本原上同调群} $H_0^k(X,\mathbb C)$，也记为 $PH^k(X)$。配对
\[
Q(\alpha,\beta)=(-1)^{k(k-1)/2}
\int_X\omega^{n-k}\wedge\alpha\wedge\beta,
\qquad\alpha,\beta\in H_0^k(X,\mathbb C),
\]
使本原上同调成为\emph{极化 Hodge 结构}。这意味着
\begin{equation}\label{eq:hodge-polarization}
Q(i^{p-q}\alpha,\overline\alpha)>0,
\qquad0\ne\alpha\in H_0^{p,q}(X),\quad p+q=k.
\end{equation}
关于极化 Hodge 结构的更多内容，见 \SourceCite{CaK,Griffiths2}。

子空间 $H^{p,q}(X)$ 在族中并不全纯地变化，但 $F^p(X)$ 的相对版本具有很好的性质。设有相对维数为 $n$ 的光滑态射 $\pi:\mathcal X\to S$，其中 $S$ 光滑且拟射影。对任意 $t\in S$，置 $X_t=\pi^{-1}(t)$。各上同调群 $H^k(X_t,\mathbb C)$ 拼成 $S$ 上的局部自由层 $\mathcal F^0=R^k\pi_*\mathbb C\otimes\mathcal O_S$。各空间 $F^p(X_t)$ 也拼成局部自由子层 $\mathcal F^p$。

由构造，$\mathcal F^0$ 含有局部系统，即局部常值层 $R^k\pi_*\mathbb C$。它唯一确定 $\mathcal F^0$ 上的平坦联络 $\nabla$，称为 \emph{Gauss--Manin 联络}；其平坦截面，也称水平截面，恰好是局部系统 $R^k\pi_*\mathbb C$ 的截面。具体地，$\nabla:\mathcal F^0\to\mathcal F^0\otimes\Omega_S^1$ 定义为
\[
\nabla(s\otimes f)=s\otimes df,
\]
其中 $s$ 是 $R^k\pi_*\mathbb C$ 的截面，$f$ 是 $S$ 上的函数。Gauss--Manin 联络满足 \emph{Griffiths 横截性}
$\nabla(\mathcal F^p)\subset\mathcal F^{p-1}\otimes\Omega_S^1$。

若置 $\mathcal H=\mathcal F^0$，则 $\mathcal H$ 有局部常值子层 $\mathcal H_{\mathbb C}=R^k\pi_*\mathbb C$；它又有整数截面构成的子层 $\mathcal H_{\mathbb Z}$，即 $R^k\pi_*\mathbb Z\to R^k\pi_*\mathbb C$ 的像。称 $(\mathcal H,\nabla,\mathcal H_{\mathbb Z},\mathcal F^\bullet)$ 为 \emph{Hodge 结构变分}。\index{Hodge结构@Hodge 结构|hyperpage}\index{Hodge结构变分@Hodge 结构变分|hyperpage}\index{Gauss--Manin联络@Gauss--Manin 联络|hyperpage}\index{Griffiths横截性@Griffiths 横截性|hyperpage}

\subsection{Picard--Fuchs 方程}\label{subsec:picard-fuchs}

本节余下部分研究中维上同调群 $H^n(X_t,\mathbb C)$ 上的 Hodge 结构变分，其中 $n=\dim X_t$。固定 $p\in S$ 处 $\mathcal F^n$ 的局部截面 $\Omega(t)$，令 $\mathcal D$ 为 $S$ 上的线性微分算子层。在 $p$ 处的局部解析坐标 $z_1,\ldots,z_r$ 下，微分算子环为
\[
\mathcal D=\mathbb C\{z_1,\ldots,z_r\}
\left[\frac\partial{\partial z_1},\ldots,
\frac\partial{\partial z_r}\right],
\]
其中 $\mathbb C\{z_1,\ldots,z_r\}$ 是 $z_1,\ldots,z_r$ 的收敛幂级数环。Gauss--Manin 联络 $\nabla$ 确定 $\mathcal O_S$ 同态 $\phi:\mathcal D\to\mathcal F^0$，对 $S$ 上的向量场 $X_i$，其规则为
\[
\phi(X_1\cdots X_\ell)
=\nabla_{X_1}\cdots\nabla_{X_\ell}\Omega(t).
\]
联络同时赋予 $\mathcal F^0$ 一个 $\mathcal D$ 模结构。

左理想 $I=\ker\phi$ 由零化 $\Omega(t)$ 的微分算子组成，称为 \emph{Picard--Fuchs 理想}。此理想在理论中起重要作用。例如，后面会看到，知道 $I$ 就能计算 Yukawa 耦合。若上述映射 $\phi$ 满射，\SourceCite{BG} 证明这发生在 Calabi--Yau 三维簇上，
\SourcePage{98}
则 $\mathcal F^0$ 的 $\mathcal D$ 模结构完全由 Picard--Fuchs 理想 $I$ 确定。

对 $D\in I$，方程 $D\cdot y=0$ 称为 \emph{Picard--Fuchs 方程}；若 $\mathcal O_S$ 的局部截面 $y$ 对所有 $D\in I$ 都满足 $D\cdot y=0$，则称其为一个\emph{解}。由于对局部常值的同调 $n$ 圈 $\gamma_t$ 有
\begin{equation}\label{eq:period-derivative}
\frac\partial{\partial z_i}\int_{\gamma_t}\Omega(t)
=\int_{\gamma_t}\nabla_{\partial/\partial z_i}\Omega(t),
\end{equation}
所以 $\Omega$ 的周期 $y=\int_{\gamma_t}\Omega(t)$ 是 Picard--Fuchs 方程的解。我们常将周期写为 $y=\langle g(t),\Omega(t)\rangle$，其中 $g(t)\in H^n(X_t,\mathbb C)$ 是 $\gamma_t$ 的 Poincaré 对偶，因此 $g(t)$ 是 $\nabla$ 的平坦截面。注意，微分算子 $D$ 属于 Picard--Fuchs 理想，当且仅当对所有周期 $y=\langle g,\Omega\rangle$ 都有 $D\cdot y=0$。

虽然这里只在局部构造了 Picard--Fuchs 方程，但若限于代数微分形式的周期，通过解析延拓，它们可作为整个 $S$ 上的代数微分方程来理解。

\subsection{Hodge 结构的退化}\label{subsec:hodge-degenerations}

设光滑族 $\pi:\mathcal X\to S$ 可补成平坦族 $\overline\pi:\overline{\mathcal X}\to\overline S$，其中 $\overline S$ 是 $S$ 的光滑紧化，边界除子 $D=\bigcup_j D_j=\overline S-S$ 具有正规交叉。

$S$ 上的丛 $\mathcal F^0$ 在 $\overline S$ 上有一个\emph{典范延拓} $\overline{\mathcal F}^0$，见 \SourceCite{Deligne1}。Gauss--Manin 联络不一定延拓为 $\overline{\mathcal F}^0$ 上的通常联络，因为它可能产生奇点。这些奇点很温和，是\IndexLocation{indexloc:1225:98}正则奇点；也就是说，$\nabla$ 延拓为映射
$\overline\nabla:\overline{\mathcal F}^0\to\overline{\mathcal F}^0\otimes\Omega_{\overline S}^1(\log D)$。
层 $\Omega_{\overline S}^1(\log D)$ 是 $\overline S$ 上有理一形式层的子层。在 $\overline S$ 的局部坐标 $z_1,\ldots,z_r$ 中，若 $D$ 由 $z_1\cdots z_\ell=0$ 定义，则 $\Omega_{\overline S}^1(\log D)$ 由
\[
\frac{dz_1}{z_1},\ldots,\frac{dz_\ell}{z_\ell},
dz_{\ell+1},\ldots,dz_r
\]
生成。

还需考虑单值性。设 $\gamma_j=\gamma_j(u)$ 是围绕边界除子 $D_j$ 的一条小回路，以 $\gamma_j(0)=\gamma_j(1)=t\in S$ 为基点。上同调类 $\eta\in H^n(X_t,\mathbb C)$ 可以唯一提升为 $[0,1]$ 上的 $\nabla$ 平坦截面 $\eta(u)\in H^n(X_{\gamma_j(u)},\mathbb C)$，满足 $\eta(0)=\eta$。\emph{\IndexLocation{indexloc:0901:98}\IndexLocation{indexloc:0902:98}\IndexLocation{indexloc:0904:98}\IndexLocation{indexloc:0905:98}单值性变换}，也称 \emph{Picard--Lefschetz 变换}，
\[
\mathcal T_j:H^n(X_t,\mathbb C)\longrightarrow H^n(X_t,\mathbb C)
\]
定义为 $\mathcal T_j(\eta)=\eta(1)$。在共轭意义下，它只依赖于 $D_j$，而不依赖于 $\gamma_j$ 的选择。\IndexLocation{indexloc:0900:98}单值性定理 \SourceCite{Landman} 表明，对某个整数 $m>0$，有
\[
(\mathcal T_j^m-I)^{n+1}=0.
\]
因此 $\mathcal T_j$ 是拟幺幂的，幺幂指数至多为 $n+1$。特别地，$\mathcal T_j$ 的特征值都是单位根。\index{单值性|hyperpage}

\subsection{极限混合 Hodge 结构}\label{subsec:limiting-hodge}

假设单值性幺幂，且在点 $p\in\overline S$ 处，$D=\overline S-S$ 由 $z_1\cdots z_r=0$ 定义；$z_1,\ldots,z_r$ 仍为 $p$ 处的局部坐标。因此可设 $S=(\Delta^*)^r$、$\overline S=\Delta^r$、$p=0$。令 $N_j=\log\mathcal T_j$，其中 $\mathcal T_j$ 为逆时针绕 $(\Delta^*)^r$ 的第 $j$ 个因子一周所得的单值性。在这种情形，$\mathcal F^0$ 的典范延拓可描述如下。设 $s$ 是 $(\Delta^*)^r$ 上 $\mathcal H=\mathcal F^0$ 的平坦多值截面，则
\[
\exp\left(-\frac1{2\pi i}\sum_j\log(z_j)N_j\right)s
\]
\SourcePage{99}
是单值的，并延拓为典范延拓 $\overline{\mathcal F}^0$ 的截面。容易看出，这个性质刻画了 $\overline{\mathcal F}^0$。

此外，幂零轨道定理 \SourceCite{Schmid} 意味着丛 $\mathcal F^p$ 延拓为 $\overline{\mathcal F}^p$，它们是典范延拓的子丛。在 $0\in\Delta^r$，$\overline{\mathcal F}^\bullet$ 给出极限 Hodge 滤过 $\overline{\mathcal F}_0^\bullet=F_{\mathrm{lim}}^\bullet$。可以证明 $N_j(F_{\mathrm{lim}}^p)\subset F_{\mathrm{lim}}^{p-1}$，因此得到线性映射
\[
N_j:F_{\mathrm{lim}}^p/F_{\mathrm{lim}}^{p+1}
\longrightarrow F_{\mathrm{lim}}^{p-1}/F_{\mathrm{lim}}^p.
\]
若 $\delta_j=z_j\partial/\partial z_j$，则由 Griffiths 横截性，
$\overline\nabla_{\delta_j}(\overline{\mathcal F}^p)\subset\overline{\mathcal F}^{p-1}$。于是 $\overline\nabla_{\delta_j}$ 诱导线性映射
\[
\overline\nabla_{\delta_j}:F_{\mathrm{lim}}^p/F_{\mathrm{lim}}^{p+1}
\longrightarrow F_{\mathrm{lim}}^{p-1}/F_{\mathrm{lim}}^p.
\]
利用幂零轨道定理，可见这两个映射满足关系
\begin{equation}\label{eq:connection-monodromy-residue}
\overline\nabla_{\delta_j}=-\frac1{2\pi i}N_j.
\end{equation}

在 $0\in\Delta^r$ 上还有自然的整数结构。设 $s$ 是 $(\Delta^*)^r$ 上 $\mathcal H_{\mathbb Z}$ 的平坦多值截面。如上，
\[
\widetilde s=\exp\left(-\frac1{2\pi i}\sum_j\log(z_j)N_j\right)s
\]
是典范延拓 $\overline{\mathcal F}^0$ 的截面，$\widetilde s(0)$ 是 $0$ 上的整元素。若用 $\exp(\tfrac1{2\pi i}\sum_j\log(z_j)N_j)$ 将其传回 $H^n(X_t,\mathbb C)$，就得到 $H^n(X_t,\mathbb Z)$ 所确定的通常整数结构。

另一个重要部分是\emph{单值性权滤过}
\[
W_\bullet:\quad W_0\subset W_1\subset\cdots\subset W_{2n-1}
\subset W_{2n}=H^n(X_t,\mathbb C),
\]
它由 $N=\sum_j a_jN_j$、$a_j>0$ 在 $H^n(X_t,\mathbb C)$ 上的作用定义。例如，
\begin{equation}\label{eq:monodromy-weight-filtration}
\begin{aligned}
W_0&=\operatorname{im}(N^n),\\
W_1&=\operatorname{im}(N^{n-1})\cap\ker N,\\
W_2&=\operatorname{im}(N^{n-2})\cap\ker N
+\operatorname{im}(N^{n-1})\cap\ker(N^2),\\
&\ \vdots\\
W_{2n-1}&=\ker(N^n).
\end{aligned}
\end{equation}
其他 $W_k$ 的公式可在 \SourceCite{Griffiths2} 中找到。单值性权滤过的主要性质为：
\begin{itemize}
\item $N(W_k)\subset W_{k-2}$；
\item $N^k$ 诱导同构
$N^k:W_{n+k}/W_{n+k-1}\simeq W_{n-k}/W_{n-k-1}$；
\item $F_{\mathrm{lim}}^\bullet$ 在 $W_k/W_{k-1}$ 上诱导权为 $k$ 的 Hodge 结构。
\end{itemize}
最后一项是说，$(W_\bullet,F_{\mathrm{lim}}^\bullet)$ 是\emph{混合 Hodge 结构}。更完整的讨论见 \SourceCite{Griffiths2}。单值性权滤过在\hyperref[sec:maximal-unipotent]{第 5.2 节}中将起重要作用。\index{极限混合Hodge结构@极限混合 Hodge 结构|hyperpage}\index{单值性权滤过|hyperpage}

\subsection{单值性与 Picard--Fuchs 方程}\label{subsec:monodromy-pf}

为简单起见，设 $\dim S=1$，令 $z$ 为以 $p\in\overline S-S$ 为中心的圆盘 $\Delta\subset\overline S$ 上的坐标。选择 $\overline{\mathcal F}^0$ 在 $\Delta$ 上的一组基 $\omega_1,\ldots,\omega_r$，则联络 $\overline\nabla$ 完全由其\emph{\IndexLocation{indexloc:0174:99}联络矩阵} $(\Gamma_{ij})$ 确定，其中 $\overline\nabla_{d/dz}\omega_i=\sum_j\Gamma_{ij}\omega_j$。

\SourcePage{100}
因为 $\overline\nabla$ 有正则奇点，$\Gamma_{ij}$ 在 $z=0$ 处至多有一阶极点。由此得到留数矩阵
\begin{equation}\label{eq:connection-residue}
\operatorname{Res}(\nabla)=\operatorname{Res}_{z=0}(\Gamma_{ij}),
\end{equation}
根据 \SourceCite{Deligne1}，它具有以下性质：
\begin{itemize}
\item $\operatorname{Res}(\nabla)$ 的特征值 $\lambda$ 是有理数，满足 $0\leq\lambda<1$；
\item $\exp(-2\pi i\operatorname{Res}(\nabla))$ 与单值性 $\mathcal T$ 共轭；
\item $\mathcal T$ 幺幂，当且仅当 $\operatorname{Res}(\nabla)$ 幂零。
\end{itemize}
实际计算时，可能不知道 $\overline{\mathcal F}^0$ 的一组基。改为假定 $\omega_1,\ldots,\omega_r$ 是去心圆盘 $\Delta^*$ 上的一组基。若联络矩阵在 $z=0$ 处至多有一阶极点，留数矩阵 $\operatorname{Res}(\nabla)$ 仍给出有用的信息：
\begin{itemize}
\item 当 $\lambda$ 遍历 $\operatorname{Res}(\nabla)$ 的特征值时，单值性 $\mathcal T$ 的特征值为 $\exp(-2\pi i\lambda)$；
\item $\mathcal T$ 幺幂，当且仅当 $\operatorname{Res}(\nabla)$ 的特征值全为整数；
\item 若 $\operatorname{Res}(\nabla)$ 的任意两个不同特征值之差都不是整数，则 $\mathcal T$ 与 $\exp(-2\pi i\operatorname{Res}(\nabla))$ 共轭。
\end{itemize}

为看出此时 Picard--Fuchs 方程的形式，令 $\delta=z\,d/dz$。对 $\overline{\mathcal F}^n$ 的固定截面 $\Omega(z)$，设在 $\overline{\mathcal F}^0$ 中有关系
\[
\nabla_\delta^m\Omega+f_1(z)\nabla_\delta^{m-1}\Omega
+\cdots+f_m(z)\Omega=0,
\]
其中各 $f_i(z)$ 都在 $z=0$ 处解析。如前所述，这意味着 $\delta^m+f_1(z)\delta^{m-1}+\cdots+f_m(z)$ 属于 Picard--Fuchs 理想，而周期 $y=\int_{\gamma_z}\Omega(z)$ 是 Picard--Fuchs 方程
\begin{equation}\label{eq:regular-pf-equation}
\delta^my+f_1(z)\delta^{m-1}y+\cdots+f_m(z)y=0
\end{equation}
的解；由于单值性对 $\gamma_z$ 的影响，这些解可能是多值的。

现在像 \SourceCite{CL} 那样，将它与常微分方程联系起来。按照微分方程标准教材中的定义，式~\eqref{eq:regular-pf-equation} 在这里具有正则奇点。其解由\emph{指标方程}控制；这个方程通过将式~\eqref{eq:regular-pf-equation} 中的 $\delta$ 替换为辅助变量 $\lambda$，并在 $z=0$ 处取系数而得到：
\[
\lambda^m+f_1(0)\lambda^{m-1}+\cdots+f_m(0)=0.
\]
设 $\lambda$ 是指标方程的一个根，而且不存在根 $\lambda'$ 使 $\lambda'-\lambda$ 为正整数。则 Picard--Fuchs 方程~\eqref{eq:regular-pf-equation} 在 $z=0$ 附近有形如 $y_0=z^\lambda f_0(z)$ 的解，$f_0$ 在 $0$ 处解析且不为零。因此，将 $y_0$ 绕 $z=0$ 解析延拓一周，得到 $\exp(2\pi i\lambda)y_0$。此外，若 $\lambda$ 是重数 $\mu>1$ 的根，则经典的 Frobenius 方法除 $y_0=z^\lambda f_0(z)$ 外，还给出其他解 $y_1,\ldots,y_{\mu-1}$；当 $1\leq i\leq\mu-1$ 时，将 $y_i$ 绕 $z=0$ 解析延拓一周，得到 $\exp(2\pi i\lambda)(y_i+y_{i-1})$。这意味着当 $z\to0$ 时，有渐近行为
\[
y_i\sim\frac c{i!}z^\lambda
\left(\frac{\log z}{2\pi i}\right)^i,
\qquad0\leq i\leq\mu-1,
\]
其中 $c$ 是非零常数。数 $\mu$ 与指标方程的友矩阵的 Jordan 标准形有关。

上述 Picard--Fuchs 方程来自 $\mathcal F^n$ 的截面 $\Omega$。若改用 $\Omega'=g(z)\Omega$，其中 $g(z)$ 在 $S$ 上亚纯，则得到新的 Picard--Fuchs 方程，
\SourcePage{101}
其解为 $\widetilde y=g(z)y$，$y$ 是式~\eqref{eq:regular-pf-equation} 的解。在这种情形，可以证明新方程仍有正则奇点，指标方程为
\[
(\lambda-\ell)^m+f_1(0)(\lambda-\ell)^{m-1}
+\cdots+f_m(0)=0,
\]
其中 $g(z)$ 在 $z=0$ 处的 Laurent 展开为 $g(z)=a_\ell z^\ell+\cdots$，$a_\ell\ne0$。

下一节将看到，在适当类型的边界点处，Picard--Fuchs 方程与单值性之间的关系尤其紧密。\index{指标方程|hyperpage}\index{Picard--Fuchs方程@Picard--Fuchs 方程|hyperpage}

\section{最大幺幂单值性}\label{sec:maximal-unipotent}

本节将发现，五次三维簇的镜像具有某种特殊的单值性行为，它对镜像对称至关重要。

\subsection{一维模}\label{subsec:one-dimensional-mum}

首先在复模空间 $S$ 为一维的 Calabi--Yau 三维簇 $V$ 的情形中讨论最大幺幂单值性。\hyperref[chap:quintic]{第 2 章}描述的五次三维簇的镜像就是一个例子。$S$ 上未必有万有族，但可在 $S$ 的有限覆盖上局部构造族。只要从 $\mathcal F^3$ 的代数截面出发，所得 Picard--Fuchs 方程就具有代数微分方程的意义，因此可解析延拓到整个 $S$。仍令 $\overline S$ 是 $S$ 的紧化，设边界点 $p\in\overline S-S$ 具有\emph{最大幺幂单值性}。这意味着单值性 $\mathcal T:H^3(V,\mathbb C)\to H^3(V,\mathbb C)$ 幺幂，而单值性定理给出 $(\mathcal T-I)^4=0$，所以“最大幺幂”意味着 $(\mathcal T-I)^3\ne0$。因此 $N=\log\mathcal T$ 也满足 $N^3\ne0$。

首要目标是解释上一节的理论如何应用于此。我们将看到，最大幺幂单值性对周期与 Picard--Fuchs 方程有一些很强的推论。首先，因为 $V$ 的复模是一维的，所以 $h^{2,1}(V)=1$，进而 $h^3(V)=4$。由于 $\mathcal T$ 最大幺幂，可找到 $H^3(V,\mathbb C)$ 的一组基 $g_0,g_1,g_2,g_3$，使 $\mathcal T$ 的矩阵为
\begin{equation}\label{eq:mum-jordan-matrix}
\mathcal T=
\begin{pmatrix}
1&1&0&0\\
0&1&1&0\\
0&0&1&1\\
0&0&0&1
\end{pmatrix}.
\end{equation}
每个 $g_i$ 都对偶于同调类 $\gamma_i\in H_3(V,\mathbb C)$。给定 $\overline{\mathcal F}^3$ 在 $z=0$ 附近的一组固定局部基 $\Omega$，希望描述周期 $y_i=\int_{\gamma_i}\Omega$ 在 $0$ 附近的行为。

先遵循 \SourceCite{Morrison2} 的方法。任意 $g\in H^3(V,\mathbb C)$ 都延拓为 $\mathcal F^0$ 的多值平坦截面 $g(z)$，而根据幂零轨道定理，$\exp(-\tfrac{\log z}{2\pi i}N)g(z)$ 在 $z=0$ 处单值且解析。$g=g_0$ 时，$\mathcal T(g_0)=g_0$ 意味着 $N(g_0)=0$，所以 $\exp(-\tfrac{\log z}{2\pi i}N)g_0(z)=g_0(z)$。因此 $g_0(z)$ 在 $z=0$ 处单值。由于
\[
y_0=\int_{\gamma_0}\Omega=\langle g_0(z),\Omega(z)\rangle,
\]
$y_0$ 在 $z=0$ 处单值且解析。后面还将看到 $y_0$ 在 $z=0$ 处不为零。再考虑 $y_1,y_2,y_3$。注意，对 $i\geq1$，有 $\mathcal T(g_i)=g_i+g_{i-1}$，且 $y_i=\langle g_i(z),\Omega(z)\rangle$，所以将 $y_i$ 绕 $z=0$
\SourcePage{102}
解析延拓一周，得到 $y_i+y_{i-1}$。此外，由于 $\exp(-\tfrac{\log z}{2\pi i}N)g_i(z)$ 单值，容易得到当 $z\to0$ 时形如
\[
y_i\sim\frac c{i!}\left(\frac{\log z}{2\pi i}\right)^i,
\qquad0\leq i\leq3,
\]
的渐近公式，其中 $c$ 是非零常数。

这使人想到\hyperref[subsec:monodromy-pf]{第 5.1.5 小节}讨论的微分方程理论。因此，下面解释此时 Picard--Fuchs 方程的形式。关键结果如下。
\begin{proposition}\label{prop:mum-indicial-equation}
固定复模空间 $S$ 为一维的 Calabi--Yau 三维簇 $V$，并在 $S$ 的一个边界点处取局部坐标 $z$，其中 Gauss--Manin 联络在该点具有正则奇点。设 $\mathcal F^3$ 的局部截面 $\Omega$ 给出 Picard--Fuchs 方程
\[
\delta^4y+f_1(z)\delta^3y+f_2(z)\delta^2y
+f_3(z)\delta y+f_4(z)y=0,
\]
其中 $\delta=z\,d/dz$，各 $f_i$ 在 $z=0$ 处解析。则：
\begin{enumerate}[label=(\roman*)]
\item 单值性 $\mathcal T$ 幺幂，当且仅当指标方程的根全为整数；
\item 单值性 $\mathcal T$ 最大幺幂，当且仅当指标方程对某个整数 $\ell$ 具有形式 $(\lambda-\ell)^4=0$。
\end{enumerate}
\end{proposition}
\begin{proof}
上述 Picard--Fuchs 方程意味着，在 $\mathcal F^0$ 中有
\begin{equation}\label{eq:mum-connection-relation}
\nabla_\delta^4\Omega+f_1(z)\nabla_\delta^3\Omega
+f_2(z)\nabla_\delta^2\Omega+f_3(z)\nabla_\delta\Omega
+f_4(z)\Omega=0.
\end{equation}
根据 \SourceCite{BG}，还知道 $\Omega,\nabla_\delta\Omega,\nabla_\delta^2\Omega,\nabla_\delta^3\Omega$ 在一般点处必定线性无关。特别地，可设它们在 $z=0$ 的去心邻域 $\Delta^*$ 中构成 $\mathcal F^0$ 的一组基；因为 $\Omega$ 未必是 $\overline{\mathcal F}^3$ 的截面，$z=0$ 处未必仍是一组基。

利用基 $\Omega,\nabla_\delta\Omega,\nabla_\delta^2\Omega,\nabla_\delta^3\Omega$，可计算 $\nabla$ 的联络矩阵。对这些基向量应用 $\nabla_{d/dz}$，并使用式~\eqref{eq:mum-connection-relation}，得到
\[
\begin{aligned}
\nabla_{d/dz}\Omega&=\frac1z\nabla_\delta\Omega,\\
\nabla_{d/dz}(\nabla_\delta\Omega)&=\frac1z\nabla_\delta^2\Omega,\\
\nabla_{d/dz}(\nabla_\delta^2\Omega)&=\frac1z\nabla_\delta^3\Omega,\\
\nabla_{d/dz}(\nabla_\delta^3\Omega)&=-\frac{f_4(z)}z\Omega
-\frac{f_3(z)}z\nabla_\delta\Omega\\
&\quad-\frac{f_2(z)}z\nabla_\delta^2\Omega
-\frac{f_1(z)}z\nabla_\delta^3\Omega.
\end{aligned}
\]
联络矩阵在 $z=0$ 处至多有一阶极点，因此可用\hyperref[subsec:monodromy-pf]{第 5.1.5 小节}发展的理论。在此基下，$\nabla$ 的\IndexLocation{indexloc:1231:102}留数矩阵为
\begin{equation}\label{eq:mum-companion-residue}
\operatorname{Res}(\nabla)=
\begin{pmatrix}
0&1&0&0\\
0&0&1&0\\
0&0&0&1\\
-f_4(0)&-f_3(0)&-f_2(0)&-f_1(0)
\end{pmatrix},
\end{equation}
它是指标方程的友矩阵。

于是，$\mathcal T$ 幺幂，当且仅当 $\operatorname{Res}(\nabla)$ 的特征值全为整数，也当且仅当指标方程的根全为整数。这证明了命题的第一部分。

\SourcePage{103}
对第二部分，设指标方程为 $(\lambda-\ell)^4=0$，其中 $\ell$ 为整数。将 $\Omega$ 替换为 $z^{-\ell}\Omega$，根据\hyperref[sec:hodge-theory]{第 5.1 节}末尾的讨论，可设指标方程为 $\lambda^4=0$。于是式~\eqref{eq:mum-companion-residue} 变为
\[
\operatorname{Res}(\nabla)=
\begin{pmatrix}
0&1&0&0\\
0&0&1&0\\
0&0&0&1\\
0&0&0&0
\end{pmatrix}.
\]
没有两个不同的根之差为整数，因此 $\operatorname{Res}(\nabla)$ 的性质表明，$\mathcal T$ 与 $\exp(-2\pi i\operatorname{Res}(\nabla))$ 共轭，从而 $\mathcal T$ 最大幺幂。

反过来，设 $\mathcal T$ 最大幺幂。只须找出指标方程为 $\lambda^4=0$ 的 Picard--Fuchs 方程。取 $\overline{\mathcal F}^3$ 的局部基 $\Omega$。我们断言，$\Omega,\nabla_\delta\Omega,\nabla_\delta^2\Omega,\nabla_\delta^3\Omega$ 在 $z=0$ 附近构成 $\overline{\mathcal F}^0$ 的一组基。

利用\hyperref[subsec:limiting-hodge]{第 5.1.4 小节}描述的单值性权滤过证明此断言。由式~\eqref{eq:monodromy-weight-filtration}，$W_6=H^3(V,\mathbb C)$，$W_5=\ker(N^3)$，其中 $N=\log\mathcal T$。式~\eqref{eq:mum-jordan-matrix} 给出 $\dim(W_6/W_5)=1$。由于 $F_{\mathrm{lim}}^\bullet$ 在 $W_6/W_5$ 上诱导权为六的 Hodge 结构，且 $F_{\mathrm{lim}}^6=F_{\mathrm{lim}}^5=F_{\mathrm{lim}}^4=0$，必有 $W_6/W_5=H^{3,3}$。$W_6/W_5$ 与 $F_{\mathrm{lim}}^3$ 均为一维，因此 $F_{\mathrm{lim}}^3\cap W_5=\{0\}$。这给出直和分解
\begin{equation}\label{eq:mum-top-splitting}
F_{\mathrm{lim}}^3\oplus W_5=H^3(V,\mathbb C).
\end{equation}

$\overline{\mathcal F}^3$ 的局部基 $\Omega$ 给出非零元素 $\Omega(0)\in F_{\mathrm{lim}}^3$。由于 $W_5=\ker(N^3)$，式~\eqref{eq:mum-top-splitting} 意味着 $N^3(\Omega(0))\ne0$；再由式~\eqref{eq:mum-jordan-matrix}，$\Omega(0),N(\Omega(0)),N^2(\Omega(0)),N^3(\Omega(0))$ 是 $H^3(V,\mathbb C)$ 的一组基。迭代式~\eqref{eq:connection-monodromy-residue} 得
\[
\nabla_\delta^j(\Omega)(0)
=\left(-\frac1{2\pi i}\right)^jN^j(\Omega(0))
\pmod{F_{\mathrm{lim}}^{4-j}}.
\]
再利用 $N(F_{\mathrm{lim}}^p)\subset F_{\mathrm{lim}}^{p-1}$ 与 $\dim F_{\mathrm{lim}}^p=4-p$，容易推出 $\Omega,\nabla_\delta\Omega,\nabla_\delta^2\Omega,\nabla_\delta^3\Omega$ 在 $z=0$ 处线性无关。因此，它们在 $z=0$ 的某个邻域构成 $\overline{\mathcal F}^0$ 的一组基。

特别地，可以用这组基表示 $\nabla_\delta^4\Omega$，得到形如式~\eqref{eq:mum-connection-relation} 的 Picard--Fuchs 方程。还需证明 $f_1(0)=f_2(0)=f_3(0)=f_4(0)=0$。这已经很容易：如上，用这组基计算 $\nabla$ 的联络矩阵与留数矩阵，得到 $\operatorname{Res}(\nabla)$ 的公式~\eqref{eq:mum-companion-residue}。因为此矩阵使用 $\overline{\mathcal F}^0$ 的一组基构造，\hyperref[sec:hodge-theory]{第 5.1 节}表明，$\mathcal T$ 幺幂意味着 $\operatorname{Res}(\nabla)$ 幂零。因此指标方程为 $\lambda^4=0$，命题得证。
\end{proof}

关于命题~\ref{prop:mum-indicial-equation} 的证明，有两点说明。首先，可以证明 $W_{2k}=W_{2k+1}$，且对 $0\leq p\leq3$ 有
\begin{equation}\label{eq:mum-hodge-tate-splitting}
F_{\mathrm{lim}}^p\oplus W_{2p-1}
=F_{\mathrm{lim}}^p\oplus W_{2p-2}=H^3(V,\mathbb C),
\end{equation}
这推广了式~\eqref{eq:mum-top-splitting}。由此可见，此时混合 Hodge 结构 $(W_\bullet,F_{\mathrm{lim}}^\bullet)$ 是 \emph{Hodge--Tate 结构}。更一般地，\SourceCite{Deligne2} 讨论了 Hodge--Tate 结构与镜像对称的关系。

第二点涉及命题~\ref{prop:mum-indicial-equation} 前的讨论：那里断言解 $y_0=\langle g_0,\Omega\rangle$ 在 $z=0$ 处不为零。为理解这一点，先注意，由式~\eqref{eq:monodromy-weight-filtration} 与
\[
\langle N(a),b\rangle+\langle a,N(b)\rangle=0,
\qquad a,b\in H^3(V,\mathbb C),
\]
容易推出 $\langle W_0,W_5\rangle=0$。

\SourcePage{104}
由于 $g_0\in W_0$，而 $\Omega(0)$ 张成 $F_{\mathrm{lim}}^3$，由式~\eqref{eq:mum-top-splitting} 以及杯积的非退化性，必有 $\langle g_0,\Omega(0)\rangle\ne0$。所以 $y_0(0)=\langle g_0,\Omega(0)\rangle$ 不为零。

举例说明命题~\ref{prop:mum-indicial-equation} 的用法：考虑\hyperref[chap:quintic]{第 2 章}的五次镜像族。由式~\eqref{eq:quintic-pf}，Picard--Fuchs 方程为
\begin{equation}\label{eq:quintic-mum-pf}
\begin{aligned}
0={}&\left(z\frac d{dz}\right)^4y
+\frac{2\cdot5^5z}{1+5^5z}\left(z\frac d{dz}\right)^3y
+\frac{7\cdot5^4z}{1+5^5z}\left(z\frac d{dz}\right)^2y\\
&+\frac{2\cdot5^4z}{1+5^5z}\left(z\frac d{dz}\right)y
+\frac{24\cdot5z}{1+5^5z}y.
\end{aligned}
\end{equation}
\hyperref[sec:hodge-examples]{第 5.4 节}将推导这个方程。这里的变量 $z$，在\hyperref[chap:quintic]{第 2 章}记为 $x$，等于 $z=\psi^{-5}$，其中 $\psi$ 是五次镜像的定义方程~\eqref{eq:quintic-pencil} 中的参数。因为指标方程为 $\lambda^4=0$，命题~\ref{prop:mum-indicial-equation} 表明 $z=0$ 处有最大幺幂单值性。

另一方面，五次镜像族还有两个边界点 $z=-5^{-5},\infty$。容易看出，$z=\infty$ 处的单值性不是最大幺幂的。在式~\eqref{eq:quintic-mum-pf} 中作坐标变换 $w=1/z$，算得指标方程的根为 $\{1/5,2/5,3/5,4/5\}$；由命题~\ref{prop:mum-indicial-equation}，无穷远点处的单值性甚至不是幺幂的。在适合 $z=-5^{-5}$ 的坐标中，指标方程的根为整数，因此单值性幺幂。不过，根并非全相等，所以根据该命题，单值性不可能最大幺幂。事实上，可以证明此时 $N^2=0$。

\subsection{$r$ 维模空间}\label{subsec:higher-dimensional-mum}

下面推广最大幺幂单值性的定义。对维数为 $d$ 的 Calabi--Yau 流形 $V$，令 $S$ 为 $V$ 的完整复模空间；$V$ 有奇点的情形将在\hyperref[chap:moduli]{第~\ref*{chap:moduli} 章}考虑。则局部形变空间 $S$ 是维数 $r=h^{d-1,1}(V)$ 的光滑流形，见 \SourceCite{Bogomolov,Tian,Todorov}。如\hyperref[sec:hodge-theory]{第 5.1 节}，令 $\overline S$ 是 $S$ 的紧化，边界除子 $D=\overline S-S$ 具有正规交叉。我们重点考察 $r=\dim S$ 个边界除子 $D_1,\ldots,D_r$ 的交点 $p\in\overline S-S$；\hyperref[chap:moduli]{第~\ref*{chap:moduli} 章}会更详细研究 $S$ 的自然紧化。我们关心以下类型的边界点。

\begin{definition}\label{def:maximal-unipotent-point}
\SourceCite{Morrison3} 若点 $p=D_1\cap\cdots\cap D_r$ 满足以下条件，则称 $p$ 为\emph{最大幺幂边界点}：
\begin{enumerate}[label=(\roman*)]
\item 单值性变换 $\mathcal T_i$ 全部幺幂；
\item 置 $N_i=\log\mathcal T_i$，令 $N=\sum_i a_iN_i$ 为所有 $a_i>0$ 的线性组合。若 $W_\bullet$ 是由 $N$ 确定的单值性权滤过~\eqref{eq:monodromy-weight-filtration}，则 $\dim W_0=\dim W_1=1$，$\dim W_2=1+r$；
\item 取 $W_2$ 的一组基 $g_0,\ldots,g_r$，使 $g_0$ 张成 $W_0$，并对 $1\leq i,j\leq r$ 用 $N_i(g_j)=m_{ij}g_0$ 定义 $m_{ij}$。则矩阵 $(m_{ij})$ 可逆。
\end{enumerate}
\end{definition}

此定义可解释如下。每个 $g_0\in W_0$ 给出 Picard--Fuchs 方程的解 $y_0=\langle g_0,\Omega\rangle$，它在最大幺幂边界点 $p$ 处解析。因此 $\dim W_0=1$ 是说，在相差常数倍的意义下，Picard--Fuchs 方程在 $p$ 处有唯一的全纯解。还应注意，每个 $g_i\in W_2$ 给出解 $\langle g_i,\Omega\rangle$，它沿各除子 $D_j$ 至多按对数增长。因此 $\dim W_2=1+r$ 意味着有 $r$ 个线性无关的、沿 $D_j$ 按对数增长的解。最后，条件 (iii) 中矩阵 $(m_{ij})$ 的可逆性使我们可以换到
\SourcePage{105}
新基
\begin{equation}\label{eq:mum-logarithmic-basis}
g_k'=\sum_{j=1}^r g_jm^{jk},\qquad
(m^{jk})=(m_{jk})^{-1},
\end{equation}
满足 $N_j(g_k')=\delta_{jk}g_0$。于是 $y_k=\langle g_k',\Omega\rangle$ 是 Picard--Fuchs 方程的解；它在 $j\ne k$ 时 $D_j$ 的一般点处全纯，而沿 $D_k$ 按对数增长。更精确地说，在 $p$ 附近可以局部写为
\[
y_k=y_0\frac{\log z_k}{2\pi i}+\text{全纯函数}.
\]
实际处理最大幺幂边界点时，寻找的正是这种情况。\hyperref[chap:moduli]{第~\ref*{chap:moduli} 章}将看到，镜像映射由函数 $y_0,y_1,\ldots,y_r$ 构造。

理解条件 (iii) 的另一种方式来自镜像对称。设 $V$ 与 $V^\circ$ 互为镜像。如\hyperref[sec:mirror-map]{第~\ref*{sec:mirror-map} 节}将看到，镜像映射把 $V^\circ$ 的 Kähler 模空间与 $V$ 的复模空间中最大幺幂边界点的一个邻域等同。在镜像一侧，\hyperref[chap:quantum-cohomology]{第~\ref*{chap:quantum-cohomology} 章}将用\IndexLocation{indexloc:1111:105}量子上同调在 $V^\circ$ 的 Kähler 模空间上构造平坦联络，称为\emph{A 模型联络}。局部地，此联络位于去心多圆盘 $(\Delta^*)^r$ 上，其坐标轴对应于 $V^\circ$ 的 Kähler 锥的一组生成元 $T_i$，$1\leq i\leq r$。每条轴产生此平坦联络的一个单值性变换 $\widetilde{\mathcal T}_i$，其对数 $\widetilde N_i=\log\widetilde{\mathcal T}_i$ 是 $\bigoplus_pH^{p,p}(V^\circ)$ 上与 $-T_i$ 作杯积的自同态。猜想镜像映射将 A 模型联络与 Gauss--Manin 联络等同，使得：
\begin{itemize}
\item $N_i$ 对应于 $\widetilde N_i=\mathord{\cdot}\cup(-T_i)$，因此 $N=\sum_i a_iN_i$ 对应于与丰沛类 $D=\sum_i a_iT_i$ 的相反数作杯积；
\item $g_0$ 对应于一个点的上同调类；
\item $g_i$ 对应于曲线 $C_i$ 的 Poincaré 对偶，其中各 $C_i$ 是 $H_2(V^\circ,\mathbb Z)$ 的一组合适的基。
\end{itemize}
于是，在镜像映射下，矩阵 $(m_{ij})$ 对应于交数矩阵 $(-T_i\cdot C_j)$；根据 Poincaré 对偶性，它必定可逆。因此，从镜像对称的角度看，定义~\ref{def:maximal-unipotent-point} 的第三个条件完全自然。

在最大幺幂单值性点处，同构 $N^d:W_{2d}/W_{2d-1}\simeq W_0$ 意味着 $N^d\ne0$；由单值性定理 $N^{d+1}=0$，所以 $N$ 是具有最大幂零指数的幂零算子。特别地，对复模为一维的三维簇，容易推出定义~\ref{def:maximal-unipotent-point} 等价于 $N^3\ne0$。更一般地，对维数为 $d$、复模为一维的 Calabi--Yau 簇，定义~\ref{def:maximal-unipotent-point} 等价于 $N^d\ne0$，见 \SourceCite{Morrison2}。

最后提到，权滤过的整数结构如何与定义~\ref{def:maximal-unipotent-point} 相联系。设在最大幺幂边界点处，$g_0,g_1,\ldots,g_r$ 是 $W_2$ 的整数格的一组 $\mathbb Z$ 基，其中 $g_0$ 生成 $W_0$ 的整数格。容易证明，由 $N_j(g_k)=m_{jk}g_0$ 定义的数 $m_{jk}$ 是整数。在 \SourceCite{Morrison6} 中，Morrison 猜想矩阵 $(m_{jk})$ 在 $\mathbb Z$ 上可逆；等价地，式~\eqref{eq:mum-logarithmic-basis} 中的 $m^{jk}$ 都是整数。这个\emph{\IndexLocation{indexloc:0637:105}\IndexLocation{indexloc:0972:105}整性猜想}已对五次三维簇的镜像得到验证 \SourceCite{Morrison2}；\SourceCite{Morrison6} 还讨论了猜想的其他情形。\index{最大幺幂单值性|hyperpage}\index{最大幺幂边界点|hyperpage}

\SourcePage{106}
\section{Griffiths--Dwork 方法}\label{sec:griffiths-dwork}

现在转向 Griffiths 与 Dwork 提出的计算 Picard--Fuchs 理想的方法。这种方法最初用于射影超曲面，近来又推广到加权射影空间中的超曲面，乃至更一般的射影环面簇中的超曲面。

\subsection{射影超曲面}\label{subsec:projective-residues}

从次数为 $\ell$ 的光滑超曲面 $V\subset\mathbb P^n$ 开始。有留数映射
\[
\operatorname{Res}:H^n(\mathbb P^n-V)\longrightarrow H^{n-1}(V),
\]
稍后将描述它。$H^n(\mathbb P^n-V)$ 中的元素可用 $\mathbb P^n$ 上沿 $V$ 有极点的全纯 $n$ 形式表示。为此定义
\[
\Omega_0=\sum_j(-1)^j x_j\,dx_0\wedge\cdots
\wedge\widehat{dx_j}\wedge\cdots\wedge dx_n.
\]
注意，$\Omega_0$ 是 $\Omega_{\mathbb P^n}^n(n+1)$ 的截面，而这个层是平凡层。这意味着，在相差常数倍的意义下，$\Omega_0$ 是唯一的次数为 $n+1$ 的齐次全纯 $n$ 形式。

若 $V$ 由次数为 $\ell$ 的齐次方程 $f=0$ 定义，则根据 \SourceCite{Griffiths1}，$H^n(\mathbb P^n-V)$ 中的元素可用形式
\[
\frac{P\Omega_0}{f^k},\qquad\deg P=k\ell-(n+1),
\]
表示，其中 $P$ 也齐次。对 $V$ 中任意拓扑 $(n-1)$ 圈 $\gamma$，令 $T(\gamma)$ 为 $\gamma$ 上的管，即包含在 $\mathbb P^n-V$ 中、位于 $\gamma$ 上的圆周丛。留数定义为
\begin{equation}\label{eq:residue-tube}
\int_\gamma\operatorname{Res}\frac{P\Omega_0}{f^k}
=\int_{T(\gamma)}\frac{P\Omega_0}{f^k}.
\end{equation}
向 $P\Omega_0/f^k$ 加上恰当形式，不改变式~\eqref{eq:residue-tube} 的右边，因此 $\operatorname{Res}$ 在上同调类上良定义。由于 $\ell H\sim V$，其中 $H$ 为超平面类，故 $\operatorname{Res}(P\Omega_0/f^k)\cdot H=0$。定义本原上同调
\[
PH^{n-1}(V)=\{\eta\in H^{n-1}(V):\eta\cdot H=0\}.
\]
若 $V$ 的维数为奇数，本原上同调等于全部中维上同调。一般而言，$\operatorname{Res}$ 的像包含在 $PH^{n-1}(V)$ 中；根据 \SourceCite{Griffiths1}，映射 $\operatorname{Res}:H^n(\mathbb P^n-V)\to PH^{n-1}(V)$ 满射。下文将遵循通常做法，常在记号中省去“$\operatorname{Res}$”，将 $P\Omega_0/f^k$ 视为 $V$ 上的上同调类或微分形式。

根据 \SourceCite{Griffiths1}，$P\Omega_0/f^k$ 实际位于 $F^{n-k}PH^{n-1}(V)$ 中。若让定义多项式 $f$ 与分子 $P$ 依赖参数 $s$ 来得到族 $V$，则计算得
\[
\nabla_{d/ds}\left(\frac{P\Omega_0}{f^k}\right)
=\frac{(-kPf'+fP')\Omega_0}{f^{k+1}}
\in\mathcal F^{n-k-1},
\]
其中撇号表示关于 $s$ 求导。这验证了射影超曲面的 Griffiths 横截性。

\SourcePage{107}
\SourceCite{Griffiths1} 给出了一个有效过程，用于模去恰当形式进行计算。对每个 $i$，令 $G_i$ 为次数 $(k-1)\ell-n$ 的齐次多项式，置
\[
\eta=\frac1{f^{k-1}}\sum_{i<j}(-1)^{i+j}(x_iG_j-x_jG_i)
\,dx_0\wedge\cdots\wedge\widehat{dx_i}\wedge\cdots
\wedge\widehat{dx_j}\wedge\cdots\wedge dx_n.
\]
则
\begin{equation}\label{eq:griffiths-pole-reduction}
d\eta=
\frac{(k-1)\bigl(\sum_jG_j\frac{\partial f}{\partial x_j}\bigr)\Omega_0}{f^k}
-\frac{\bigl(\sum_j\frac{\partial G_j}{\partial x_j}\bigr)\Omega_0}{f^{k-1}}.
\end{equation}
关键在于：利用式~\eqref{eq:griffiths-pole-reduction}，当 $P$ 位于 Jacobian 理想
\[
J(f)=\left(\frac{\partial f}{\partial x_0},\ldots,
\frac{\partial f}{\partial x_n}\right)
\]
中时，任何形式 $P\Omega_0/f^k$ 都可模去恰当形式，约化为极点阶数较低的形式。反过来也成立：若 $P\Omega_0/f^k$ 可模去恰当形式，约化为 $F^{n-k+1}$ 中的元素，则 $P$ 位于 Jacobian 理想。换言之，映射 $P\mapsto P\Omega_0/f^k$ 在相应的分次商上定义同构
\begin{equation}\label{eq:jacobian-hodge-isomorphism}
(S/J(f))_{k\ell-(n+1)}\simeq PH^{n-k,k-1}(V),
\end{equation}
其中 $S=\mathbb C[x_0,\ldots,x_n]$，下标表示次数为 $k\ell-(n+1)$ 的分次部分。因此，选择 $(S/J(f))_{k\ell-(n+1)}$ 的一组基之后，任何 $\deg P=k\ell-(n+1)$ 的形式 $P\Omega_0/f^k$，都可以模去恰当形式与极点阶数较低的形式，约化为来自这组基的形式的线性组合。

下面计算 Picard--Fuchs 方程。为简单起见，在 $V$ 为 $\mathbb P^n$ 中 Calabi--Yau 超曲面的情形说明方法，这意味着 $\ell=n+1$。此时 $\omega=\operatorname{Res}(\Omega_0/f)$ 是 $V$ 上的全纯 $(n-1)$ 形式。\emph{Griffiths--Dwork 方法}按如下步骤进行：
\begin{itemize}
\item 设 $V$ 由依赖参数 $s\in S$ 的单个方程 $f$ 定义，其中 $S$ 是一维的。当 $s$ 变化时，$\omega=\operatorname{Res}(\Omega_0/f)$ 是 $f$ 所定义的超曲面族 $V$ 上的全纯 $(n-1)$ 形式。假定 $f$ 对 $s$ 的依赖为多项式。
\item 为 $V$ 的本原上同调选择一组基，用一组便于计算的形式 $\omega_i=P_i\Omega_0/f^{k_i}$ 表示，$1\leq i\leq r$。这里 $r=\dim PH^{n-1}(V)$；$V$ 的维数为偶数时 $r=h^{n-1}(V)-1$，为奇数时 $r=h^{n-1}(V)$。
\item 反复对 $\omega$ 关于 $s$ 求导，得到 $\mathcal F^0$ 的截面
\[
\omega,\nabla_{d/ds}\omega,\ldots,\nabla_{d/ds}^r\omega.
\]
这些截面都可模去恰当形式，用选定的基表示。方法如下：从最高阶极点向下处理，在模 $J(f)$ 的意义下，寻找 $\nabla_{d/ds}^i\omega$ 的分子多项式与选定基的分子之间的关系；最方便的是在环 $\mathbb C(s)[x_0,\ldots,x_n]$ 中使用 Gröbner 基方法。再用式~\eqref{eq:griffiths-pole-reduction}，将求导所得的形式模去恰当形式，表示为基元素的线性组合加上极点阶数较低的形式 $\eta$。线性组合的系数是参数 $s$ 的有理函数。然后对 $\eta$ 重复这个过程。
\SourcePage{108}
\item 由于有 $r+1$ 个截面，而 $\mathcal F^0$ 的秩为 $r$，$\omega,\nabla_{d/ds}\omega,\ldots,\nabla_{d/ds}^r\omega$ 之间必定存在一个系数属于 $\mathbb C(s)$ 的线性关系。这个关系给出 Picard--Fuchs 理想中的一个元素，所得常微分方程就是 Picard--Fuchs 方程。
\end{itemize}
下一节将把 Griffiths--Dwork 方法用于五次三维簇的镜像。此外，当超曲面的定义方程 $f$ 依赖多个参数时，容易调整这个方法，得到 Picard--Fuchs 理想中的元素。

遵循 \SourceCite{Dolgachev2}，不难把 Griffiths--Dwork 方法推广到加权射影空间中的超曲面。若加权射影空间为 $\mathbb P(q_0,\ldots,q_n)$，有
\[
\Omega_0=\sum_{j=0}^n(-1)^{j-1}q_jx_j\,dx_0\wedge\cdots
\wedge\widehat{dx_j}\wedge\cdots\wedge dx_n
\]
其余修改都很直接。下一节将给出说明这个方法的例子。
\index{Griffiths--Dwork方法@Griffiths--Dwork 方法|hyperpage}\index{本原上同调|hyperpage}\index{留数映射|hyperpage}

\subsection{环面超曲面}\label{subsec:toric-residues}

现在说明 Griffiths--Dwork 方法如何应用于环面簇中的丰沛超曲面。我们遵循 \SourceCite{BC} 的处理。设 $X=X_\Sigma$ 是一个 $n$ 维射影单纯环面簇，记 $\Sigma(1)=\{\rho_1,\ldots,\rho_r\}$，$S=\mathbb C[x_1,\ldots,x_r]$ 为~\hyperref[subsec:homogeneous-ring]{第 3.2.3 小节}定义的 $X$ 的齐次坐标环。取 $\beta\in A_{n-1}(X)$，并假定 $\beta$ 是 $X$ 上某个丰沛线丛的第一 Chern 类。于是 $S_\beta$ 是次数为 $\beta$ 的齐次多项式空间，这里的齐次性采用~\hyperref[subsec:homogeneous-ring]{第 3.2.3 小节}的含义：回顾 $S$ 按 $A_{n-1}(X)$ 分次。一般的 $f\in S_\beta$ 定义一个拟光滑超曲面 $V=\{f=0\}$，其除子类为 $\beta$。

Griffiths--Dwork 方法的思路可以自然地推广到这里。首先需要定义合适的形式 $\Omega_0$。可取 $\Omega_0$ 为 $\widehat\Omega_X^n$ 的次数为 $\beta_0=\sum_i[D_i]$ 的齐次截面，因为
\[
\widehat\Omega_X^n=\mathcal O_X\Bigl(-\sum_iD_i\Bigr).
\]
这里 $D_i$ 是对应于 $\rho_i$ 的环面除子。若记 $\beta_i=\deg x_i$，则 $\beta_0=\sum_i\beta_i$。

为给出 $\Omega_0$ 的显式公式，固定 $M$ 的一组基 $e_1,\ldots,e_n$。对指标集 $I=\{i_1,\ldots,i_n\}\subset\{1,\ldots,r\}$，定义
\[
\det(e_I)=\det\bigl(\langle e_j,\rho_{i_k}\rangle\bigr)_{1\leq j,k\leq n},
\]
其中也用 $\rho_i$ 表示相应射线的本原生成元。置
\[
dx_I=dx_{i_1}\wedge\cdots\wedge dx_{i_n},\qquad
\widehat x_I=\prod_{i\notin I}x_i.
\]
于是定义
\begin{equation}\label{eq:toric-volume-form}
\Omega_0=\sum_{|I|=n}\det(e_I)\widehat x_I\,dx_I.
\end{equation}
注意，$\Omega_0$ 的次数为 $\beta_0$。此外，如 \SourceCite{BC} 所解释的，这个表达式确实在 $X$ 上定义了一个带有相应扭曲的 $n$ 形式，而不只是 $\mathbb C^r$ 上的形式。于是，$X$ 上沿 $V$ 有极点的 $n$ 形式可写成
\[
\frac{P\Omega_0}{f^k},\qquad \deg P=k\beta-\beta_0.
\]

\SourcePage{109}
为降低极点阶数，对每个 $1\leq i\leq r$，需要一个次数为 $\beta_0-\beta_i$ 的 $(n-1)$ 形式 $\Omega_i$；这里不讨论其精确定义。关键是以下公式：
\begin{equation}\label{eq:toric-pole-reduction}
d\left(\frac{P\Omega_i}{f^k}\right)
=\frac{\bigl(f\frac{\partial P}{\partial x_i}
-kP\frac{\partial f}{\partial x_i}\bigr)\Omega_0}{f^{k+1}}.
\end{equation}
在这里，定义 Jacobian 理想
\[
J(f)=\left\langle\frac{\partial f}{\partial x_1},\ldots,
\frac{\partial f}{\partial x_r}\right\rangle,
\]
可以结合 Gröbner 基方法与 Jacobian 理想，降低极点阶数，从而将任意形式模去恰当形式后，用一组给定的基表示。

在环面情形中，我们采用一个稍有不同的本原上同调概念：
\[
PH^{n-1}(V)=H^{n-1}(V)\big/\operatorname{im}H^{n-1}(X).
\]
当 $V$ 为奇数维时，这就是通常的上同调。同样有留数映射
\[
\operatorname{Res}:H^n(X-V)\longrightarrow PH^{n-1}(V).
\]
在经典情形中它是同构，但在环面情形中有正合列
\[
0\longrightarrow H^{n-2}(X)\xrightarrow{\,\smile[V]\,}H^n(X)
\longrightarrow H^n(X-V)\xrightarrow{\operatorname{Res}}
PH^{n-1}(V)\longrightarrow0.
\]
由于 $H^n(X)=H^{n/2,n/2}(X)$，式~\eqref{eq:jacobian-hodge-isomorphism} 相应地变为同构
\[
(S/J(f))_{k\beta-\beta_0}\simeq PH^{n-k,k-1}(V),
\qquad k\ne n/2+1.
\]
而当 $k=n/2+1$ 时，有正合列
\[
\begin{split}
0\longrightarrow H^{n-2}(X)\longrightarrow H^n(X)
\longrightarrow (S/J(f))_{(n/2+1)\beta-\beta_0}\\
\longrightarrow PH^{n/2-1,n/2}(V)\longrightarrow0.
\end{split}
\]
（见 \SourceCite{BC}）。注意，只有在 $n$ 为偶数时才需要考虑这个正合列；此时超曲面 $V$ 为奇数维，例如 Calabi--Yau 三维簇。

也可以稍微改变 $J(f)$，从而避免使用这个正合列。用以下理想商代替 Jacobian 理想：
\[
J_1(f)=\left\langle x_1\frac{\partial f}{\partial x_1},\ldots,
x_r\frac{\partial f}{\partial x_r}\right\rangle:(x_1\cdots x_r).
\]
根据 \SourceCite{BC}，若 $f$ 非退化，即 $V$ 拟光滑且与 $X$ 的每个轨道拟横截地相交，则对所有 $k$ 都有
\[
(S/J_1(f))_{k\beta-\beta_0}\simeq PH^{n-k,k-1}(V).
\]
为解释理想 $J_1(f)$ 的出现，注意算子 $x_i\partial/\partial x_i$ 在环面作用下不变，而取理想商是使 $J(f)\subset J_1(f)$ 成立所必需的。此外，$J_1(f)$ 也自然出现在规范线性 $\sigma$ 模型的研究中，见 \SourceCite{MP1}。

因此，使用 $J(f)$（$n$ 为奇数时）或 $J_1(f)$（任意 $n$），就能在任何射影单纯环面簇中，对丰沛超曲面应用 Griffiths--Dwork 方法来计算 Picard--Fuchs 方程。

遗憾的是，这个方法要求 $V$ 丰沛且 $X$ 单纯。例如，在 Batyrev 镜像构造的情形中，回顾~\hyperref[sec:batyrev-mirror]{第 4.1 节}：从自反多胞形 $\Delta$ 出发，$\overline V$ 是 $\mathbb P_\Delta$ 中的丰沛反典范超曲面；对 $\Delta$ 的法扇作极大射影细分 $\Sigma$，则 $V\subset X_\Sigma$ 是 $\overline V$ 的严格变换，也是一个极小 Calabi--Yau 簇。
\SourcePage{110}
虽然 $\overline V$ 丰沛，$\mathbb P_\Delta$ 却未必单纯；而虽然 $X_\Sigma$ 单纯，$V$ 却未必丰沛。所以，这个方法可能对两者都不适用。\index{Jacobian理想@Jacobian 理想|hyperpage}

\section{例题}\label{sec:hodge-examples}

现在给出两个例子，说明如何使用上一节的方法计算 Picard--Fuchs 方程。

\begin{example}\label{ex:quintic-pf-method}
首先遵循 \SourceCite{Morrison1} 的思路，考虑五次镜像族。由式~\eqref{eq:quintic-pencil}，五次镜像由方程
\begin{equation}\label{eq:quintic-pf-family}
f=x_1^5+x_2^5+x_3^5+x_4^5+x_5^5+\psi x_1x_2x_3x_4x_5=0
\end{equation}
模去式~\eqref{eq:greene-plesser-group} 中给出的群 $G$ 的作用而得到。更精确地说，如\hyperref[sec:quintic-revisited]{第 4.2 节}所述，若 $W\subset\mathbb P^4$ 是式~\eqref{eq:quintic-pf-family} 定义的超曲面，则 $V^\circ$ 是
$W/G\subset\mathbb P^4/G=\mathbb P_{\Delta^\circ}$ 的消解，其中 $\Delta^\circ$ 的顶点由式~\eqref{eq:quintic-polar-vertices} 给出。再由\hyperref[chap:mirror-constructions]{第~\ref*{chap:mirror-constructions} 章}的式~\eqref{eq:batyrev-hn21}，有
\begin{equation}\label{eq:example-polynomial-hodge}
\begin{split}
h^{2,1}(V^\circ)
&=l(\Delta^\circ)-5-\sum_{\Gamma^\circ}l^*(\Gamma^\circ)
+\sum_{\Theta^\circ}l^*(\Theta^\circ)l^*(\widehat\Theta^\circ)\\
&=h_{\mathrm{poly}}^{2,1}(V^\circ)
+\sum_{\Theta^\circ}l^*(\Theta^\circ)l^*(\widehat\Theta^\circ).
\end{split}
\end{equation}
这里 $\Gamma^\circ$ 遍历 $\Delta^\circ$ 的余维 $1$ 的面，$\Theta^\circ$ 遍历其余维 $2$ 的面，$\widehat\Theta^\circ$ 是 $\Delta$ 中与 $\Theta^\circ$ 对偶的一维面。容易算出，$\Delta^\circ$ 的整点只有式~\eqref{eq:quintic-polar-vertices} 所列的五个顶点和原点 $(0,0,0,0)$，共六个，因而
\[
h^{2,1}(V^\circ)=h_{\mathrm{poly}}^{2,1}(V^\circ)=6-5-0+0=1.
\]
所以 $V^\circ$ 的复模空间为一维，来自式~\eqref{eq:quintic-pf-family} 中参数 $\psi$ 的变化。回顾\hyperref[chap:quintic]{第~\ref*{chap:quintic} 章}，$z=\psi^{-5}$ 是模空间的一个局部参数；\hyperref[chap:quintic]{第~\ref*{chap:quintic} 章}中记作 $x=\psi^{-5}$ 的变量，现在改记为 $z$。

因为 $V^\circ$ 是 $\mathbb P_{\Delta^\circ}$ 的一个环面爆破中的非丰沛除子，无法对它应用 Griffiths--Dwork 方法。幸而 $W/G$ 丰沛，且 $\mathbb P_{\Delta^\circ}$ 单纯，所以改取 $V^\circ=W/G$，因为消解不改变 Picard--Fuchs 方程。这样便可应用\hyperref[sec:griffiths-dwork]{第 5.3 节}的方法，只须用 $\mathbb P^4-W$ 上的 $G$ 不变形式 $P\Omega_0/f^k$ 表示 $V^\circ$ 上的上同调类。

从 $G$ 不变三形式 $\omega_1=\psi\Omega_0/f$ 开始，它已在\hyperref[chap:quintic]{第 2 章}的式~\eqref{eq:quintic-residue} 中出现。目标是找出
$\omega_1,\delta\omega_1,\ldots,\delta^4\omega_1$ 之间的关系，其中
$\delta=z\,d/dz=-\frac15\psi\,d/d\psi$。为此，需要一组上同调的基。考虑以下 $G$ 不变形式：
\[
\omega_j=(-1)^{j-1}\psi^j(x_1x_2x_3x_4x_5)^{j-1}
\frac{\Omega_0}{f^j},\qquad j=1,\ldots,4.
\]
选择因子 $\psi^j$，是为了使 $\omega_j$ 在~\hyperref[sec:quintic-mirror]{第 2.2 节}所述作用
$x_1\mapsto\mu^{-1}x_1$、$\psi\mapsto\mu\psi$（$\mu^5=1$）下不变。这个变换诱导同构 $V^\circ_\psi\simeq V^\circ_{\mu\psi}$；特别地，这说明 Picard--Fuchs 方程可以用 $z=\psi^{-5}$ 改写。
注意，$\operatorname{Res}(\omega_j)\in F^{4-j}H^3(V^\circ)$，一次简单的 Gröbner 基计算表明，这个类在 $H^{4-j,j-1}(V^\circ)$ 中的投影非零。由于 $h^{3,0}=h^{2,1}=h^{1,2}=h^{0,3}=1$，$\omega_j$（$j=1,2,3,4$）的留数组成 $H^3(V^\circ)$ 的一组基。

\SourcePage{111}
因此，对 $k=1,\ldots,4$，任何极点阶数为 $k$ 的 $G$ 不变形式，都可以模去恰当形式和极点阶数较低的形式，约化为 $\omega_k$ 的倍数；而具有任意阶极点的 $G$ 不变形式，都可以模去恰当形式，表示为这四个形式的线性组合。可以结合式~\eqref{eq:griffiths-pole-reduction} 中的极点约化公式与 Gröbner 基方法，完成这些计算。

现在把这个过程用于 $\omega_1,\delta\omega_1,\ldots,\delta^4\omega_1$。对任意 $j$，有
\begin{equation}\label{eq:quintic-basis-derivative}
\delta\omega_j=-\frac j5(\omega_j+\omega_{j+1}).
\end{equation}
由此立刻得到
\begin{equation}\label{eq:quintic-derivative-basis}
\begin{pmatrix}\omega_1\\\delta\omega_1\\\delta^2\omega_1\\\delta^3\omega_1\end{pmatrix}
=\begin{pmatrix}
1&0&0&0\\
-1/5&-1/5&0&0\\
1/25&3/25&2/25&0\\
-1/125&-7/125&-12/125&-6/125
\end{pmatrix}
\begin{pmatrix}\omega_1\\\omega_2\\\omega_3\\\omega_4\end{pmatrix}.
\end{equation}
这样便处理了 $\omega_1,\delta\omega_1,\delta^2\omega_1,\delta^3\omega_1$。至于 $\delta^4\omega_1$，反复进行 Gröbner 基计算并降低极点阶数，模去恰当形式后得到
\[
\delta^4\omega_1
=\frac{5z}{1+5^5z}\omega_1+\frac{75z}{1+5^5z}\omega_2
\quad+\frac{250z}{1+5^5z}\omega_3+\frac{300z}{1+5^5z}\omega_4.
\]
因此，
\[
\delta^4\omega_1
=\begin{pmatrix}
\dfrac{5z}{1+5^5z}&\dfrac{75z}{1+5^5z}&
\dfrac{250z}{1+5^5z}&\dfrac{300z}{1+5^5z}
\end{pmatrix}
\begin{pmatrix}\omega_1\\\omega_2\\\omega_3\\\omega_4\end{pmatrix}.
\]
将式~\eqref{eq:quintic-derivative-basis} 代入，便得到 Picard--Fuchs 方程~\eqref{eq:quintic-mum-pf}。
\end{example}

\begin{example}\label{ex:weighted-pf-method}
现在考虑 $\mathbb P(1,1,2,2,2)$ 的一个爆破中的 Calabi--Yau \IndexLocation{indexloc:1020:111}超曲面的镜像。稍后将看到，镜像的复模空间是二维的，这意味着 Picard--Fuchs 方程是偏微分方程。

对于 $\mathbb P(1,1,2,2,2)$ 的标准扇，一维锥的生成元是
\begin{equation}\label{eq:weighted-pf-rays}
\{(-1,-2,-2,-2),(1,0,0,0),(0,1,0,0),(0,0,1,0),(0,0,0,1)\}.
\end{equation}
这些点是一个\IndexLocation{indexloc:1222:111}自反多胞形 $\Delta^\circ\subset N_{\mathbb R}$ 的顶点。根据~\hyperref[sec:quintic-revisited]{第 4.2 节}的讨论，$\mathbb P_{\Delta^\circ}$ 是加权射影空间 $\mathbb P(1,1,2,2,2)$ 对某个有限群 $H$ 的商。
$\Delta^\circ$ 中只有七个整点：五个顶点、点
\[
(0,-1,-1,-1)=\tfrac12(-1,-2,-2,-2)+\tfrac12(1,0,0,0)
\]
以及原点。由式~\eqref{eq:example-polynomial-hodge}，可得
$h^{2,1}(V^\circ)=h_{\mathrm{poly}}^{2,1}(V^\circ)=2$。所以 $V^\circ$ 有二维复模；而且这些模是多项式模，也就是说，可以通过改变 $V^\circ$ 的定义方程来检测它们。

如前所述，可以假定 $V^\circ$ 是 $\mathbb P_{\Delta^\circ}$ 中的反典范超曲面。因为它丰沛，且 $\mathbb P_{\Delta^\circ}$ 单纯，所以 Griffiths--Dwork 方法适用。

\SourcePage{112}
由于 $\Delta^\circ$ 是式~\eqref{eq:weighted-pf-rays} 中各点的凸包，计算得极多胞形 $\Delta$ 是以下各点的凸包：
\begin{equation}\label{eq:weighted-pf-polar}
\begin{split}
\{&(-1,-1,-1,-1),(7,-1,-1,-1),(-1,3,-1,-1),\\
&(-1,-1,3,-1),(-1,-1,-1,3)\}.
\end{split}
\end{equation}
按给定顺序将 $\Delta$ 的顶点记为 $v_i$。这些 $v_i$ 生成 $\mathbb P_{\Delta^\circ}$ 的扇的一维锥。注意，$\mathbb P_\Delta=\mathbb P(1,1,2,2,2)$；若 $\Sigma$ 是 $\Delta$ 的法扇的唯一极大射影细分，则 $X_\Sigma$ 是 $\mathbb P(1,1,2,2,2)$ 的一次爆破，具体由向式~\eqref{eq:weighted-pf-rays} 确定的扇中插入射线 $(0,-1,-1,-1)$ 得到。由此可知，$V^\circ$ 是本例第一句话中所描述的族。

$\mathbb P_{\Delta^\circ}$ 的齐次坐标环就是 $\mathbb C[x_1,\ldots,x_5]$，按 $A_3(\mathbb P_{\Delta^\circ})$ 分次；这里这个群含有挠元。需要确定次数
$\beta_0=\sum_{i=1}^5\deg(x_i)\in A_3(\mathbb P_{\Delta^\circ})$ 的单项式。具体而言，它们是在
$\mathbb P_{\Delta^\circ}=\mathbb P(1,1,2,2,2)/H$ 中次数为 $8$ 的 $H$ 不变单项式。最简便的办法是使用同构~\eqref{eq:toric-graded-sections}；由此得出，所求单项式为
\[
\prod_i x_i^{\langle v,v_i\rangle+1},\qquad v\in\Delta^\circ\cap N.
\]
由于已经知道 $\Delta^\circ\cap N$ 的各点，所以得到单项式
\[
x_1^8,\quad x_2^8,\quad x_3^4,\quad x_4^4,\quad x_5^4,
\quad x_1^4x_2^4,\quad x_1x_2x_3x_4x_5.
\]
它们恰好是加权射影空间中权为 $8$ 的 $H$ 不变单项式。从加权射影空间的角度看，这正是 Greene--Plesser 轨形构造。一般地，可以证明：对允许 Fermat 超曲面的加权射影空间，Batyrev 镜像总与 Greene--Plesser 构造一致（见 \SourceCite{BK2}）。

因此，$\mathbb P_{\Delta^\circ}$ 中的一般反典范超曲面由这些 $H$ 不变单项式的线性组合给出。利用环面作用缩放变量 $x_i$，可以将反典范超曲面的方程写成
\begin{equation}\label{eq:weighted-pf-family}
f=z_2x_1^8+x_2^8+z_1x_3^4+x_4^4+x_5^4
+x_1x_2x_3x_4x_5+x_1^4x_2^4=0,
\end{equation}
其中 $z_1,z_2$ 是参数。之所以这样选择，是为了使 $z_1=z_2=0$ 为\IndexLocation{indexloc:0769:112}最大幺幂边界点。记 $\delta_i=z_i\,d/dz_i$。Hodge 数为
$h^{3,0}=h^{0,3}=1$，以及前面算出的 $h^{2,1}=h^{1,2}=2$。选择以下基：
\[
\begin{aligned}
\omega_1&=\frac{\Omega_0}{f},\\
\omega_2&=\delta_1\omega_1=-\frac{z_1x_3^4\Omega_0}{f^2},\\
\omega_3&=\delta_2\omega_1=-\frac{z_2x_1^8\Omega_0}{f^2},\\
\omega_4&=\delta_1^2\omega_1=\omega_2+\frac{2z_1^2x_3^8\Omega_0}{f^3},\\
\omega_5&=\delta_1\delta_2\omega_1=\frac{2z_1z_2x_1^8x_3^4\Omega_0}{f^3},\\
\omega_6&=\delta_1^2\delta_2\omega_1
=\omega_5-\frac{6z_1^2z_2x_1^8x_3^8\Omega_0}{f^4}.
\end{aligned}
\]
将各种 $\delta_i^j\omega_1$ 用这组基表示，就能得到 Picard--Fuchs 方程。

\SourcePage{113}
例如，计算得
\[
\delta_2^2\omega_1=\omega_3+\frac{2z_2^2x_1^{16}\Omega_0}{f^3}.
\]
相应地，尝试利用 Jacobian 理想，将 $x_1^{16}$ 约化为只含 $x_3^8$ 和 $x_1^8x_3^4$ 的表达式。使用 Gröbner 基方法，得到关系
\[
z_2(1-4z_2)x_1^{16}-z_1^2x_3^8+4z_1z_2x_1^8x_3^4\in J(f).
\]
事实上，容易验证
\[
z_2(1-4z_2)x_1^{16}-z_1^2x_3^8+4z_1z_2x_1^8x_3^4
=-\frac1{16}\left(\sum_{i=1}^3A_ix_i\frac{\partial f}{\partial x_i}\right),
\]
其中
\[
\begin{aligned}
A_1&=2(4z_2-1)x_1^8+x_1x_2x_3x_4x_5-4x_1^4x_2^4,\\
A_2&=2x_1^8,\\
A_3&=4z_1x_3^4-16z_2x_1^8-x_1x_2x_3x_4x_5.
\end{aligned}
\]
将它与式~\eqref{eq:toric-pole-reduction} 结合，便能降低 $\delta_2^2\omega_1$ 的极点阶数。反复进行这个过程，继续降低极点阶数，最终得到 Picard--Fuchs 方程
\begin{equation}\label{eq:weighted-pf-second}
(1-4z_2)\delta_2^2\omega_1
-z_2(\delta_1^2-4\delta_1\delta_2-\delta_1+2\delta_2)\omega_1=0.
\end{equation}

\SourceCite{CdFKM} 用其他方法详细研究了这个例子，所用方程~\eqref{eq:weighted-pf-second} 的形式略有不同（另见 \SourceCite{HKTY1}）。如前所述，可以证明 $z_1=z_2=0$ 是一个最大幺幂边界点。为了详细验证，还需要另一个 Picard--Fuchs 方程。这里暂且只指出，从式~\eqref{eq:weighted-pf-second} 的形式很快可得 $N_2^2=0$，所以 $N_2$ 单独不能定义最大幺幂边界点的单值性权滤过。但是，任意正线性组合 $a_1N_1+a_2N_2$ 都满足
$(a_1N_1+a_2N_2)^3\ne0$，并且确实定义了所需类型的单值性权滤过。

我们将在例题~\ref{ex:weighted-gkz} 中再次讨论这个例子，并在下一章继续讨论。注意，具有正规交叉边界的模空间光滑紧化仍然只有一个最大幺幂边界点。各对边界除子的其他交点都不是最大幺幂的。
\end{example}

\section{超几何方程}\label{sec:hypergeometric}

现在描述 Gelfand、Kapranov 与 Zelevinsky 的超几何方程组，并讨论镜像对称中自然出现的几个例子。

\subsection{$\mathcal A$ 方程组}\label{subsec:a-system}

设 $\mathcal A=\{v_1,\ldots,v_k\}\subset\mathbb Z^{n+1}$ 是位于同一个整仿射超平面中的 $k>n+1$ 个点。将 $v_i$ 的坐标记为 $v_{ij}$。对每个 $v_i\in\mathcal A$ 引入变量 $\lambda_i$，并固定向量
$\widehat\beta=(\widehat\beta_1,\ldots,\widehat\beta_{n+1})\in\mathbb C^{n+1}$。这些数据确定一个关于 $\lambda_i$ 的微分方程组，称为\emph{超几何方程的 $\mathcal A$ 方程组}。用 $\partial_i$ 表示微分算子 $\partial/\partial\lambda_i$，用 $\delta_i$ 表示对数导数 $\lambda_i\partial_i=\lambda_i\partial/\partial\lambda_i$。

$\mathcal A$ 方程组由两种微分算子构成，记为 $Z_j$ 与 $\Box_\ell$。算子 $Z_j$ 定义为
\begin{equation}\label{eq:gkz-euler-operators}
Z_j=\left(\sum_i v_{ij}\delta_i\right)-\widehat\beta_j,
\qquad j=1,\ldots,n+1.
\end{equation}

\SourcePage{114}
为描述 $\Box_\ell$，还需要一些记号。令 $\Lambda$ 为 $\mathcal A$ 中各元素的关系格：
\begin{equation}\label{eq:gkz-relation-lattice}
\Lambda=\left\{\ell=(\ell_i)\in\mathbb Z^k:\sum_i\ell_iv_i=0\right\}.
\end{equation}
根据关于 $\mathcal A$ 的假设，$\Lambda\ne\{0\}$。然后定义
\begin{equation}\label{eq:gkz-box-operators}
\Box_\ell=\prod_{\ell_i>0}\partial_i^{\ell_i}
-\prod_{\ell_i<0}\partial_i^{-\ell_i},\qquad\ell\in\Lambda.
\end{equation}

\begin{definition}\label{def:gkz-system}
\emph{超几何方程的 $\mathcal A$ 方程组}是微分方程组
$Z_j\Phi=0$（$1\leq j\leq n+1$）与 $\Box_\ell\Phi=0$（$\ell\in\Lambda$），其中 $Z_j$ 与 $\Box_\ell$ 分别由式~\eqref{eq:gkz-euler-operators} 与式~\eqref{eq:gkz-box-operators} 定义。
\end{definition}
根据 Gelfand、Kapranov 与 Zelevinsky 的工作 \SourceCite{GKZ1}，这个方程组通常也称为 \emph{GKZ 方程组}。\index{超几何方程|hyperpage}\index{GKZ方程组@GKZ 方程组|hyperpage}

\hypertarget{source:114-batch3}{}
\subsection{Calabi--Yau 环面超曲面}\label{subsec:toric-gkz}

现在设 $\Delta\subset M_{\mathbb R}$ 是一个 $n$ 维自反多胞形，$X_\Sigma$ 是对应于 $\Delta$ 的法扇的一个极大射影细分 $\Sigma$ 的环面簇。考虑 $X_\Sigma$ 中的反典范超曲面族；由命题~\ref{prop:toric-cy-construction}，这些超曲面是极小 Calabi--Yau 子簇。我们的任务是为这个族求出 Picard--Fuchs 方程。如前所述，这些方程与细分的选择无关，而且与 $\mathbb P_\Delta$ 中反典范超曲面族 $V$ 的方程相同；由命题~\ref{prop:toric-cy-construction}，后者也是 Calabi--Yau 簇族。因此，只需研究 $\mathbb P_\Delta$ 中的这个 Calabi--Yau 超曲面族。

将 $\Delta\cap M$ 中的点 $m_i$ 排序，使 $m_0=0$。由于 $\Delta$ 自反，反典范丛的截面可以与以下 Laurent 多项式等同：
\begin{equation}\label{eq:toric-gkz-laurent}
f=\sum_i\lambda_it^{m_i}=\lambda_0+\sum_{i>0}\lambda_it^{m_i}\in L(\Delta),
\end{equation}
其中 $\lambda_i$ 是参数。于是，Calabi--Yau 簇族 $V\subset\mathbb P_\Delta$ 由方程 $f=0$ 定义，其中 $f\in L(\Delta)$。

为了将它与超几何方程联系起来，令
\begin{equation}\label{eq:toric-gkz-configuration}
\mathcal A=(\Delta\cap M)\times\{1\}\subset M\times\mathbb Z\simeq\mathbb Z^{n+1},
\end{equation}
并置 $\widehat\beta=(0,\ldots,0,-1)$。注意，$\ell\in\Lambda$ 同时意味着 $\sum_i\ell_im_i=0$ 与 $\sum_i\ell_i=0$。后一个等式说明，式~\eqref{eq:gkz-box-operators} 中的 $\Box_\ell$ 是齐次微分算子。

我们断言，$\mathcal A$ 方程组中的方程是 Picard--Fuchs 方程；也就是说，$V$ 上某些微分形式的周期满足超几何方程的 $\mathcal A$ 方程组。

遵循 \SourceCite{Batyrev3}，使用 $n$ 形式
\begin{equation}\label{eq:toric-gkz-rational-form}
\omega=\frac1f\eta,\qquad
\eta=\frac{dt_1}{t_1}\wedge\cdots\wedge\frac{dt_n}{t_n},
\end{equation}
其中 $(t_1,\ldots,t_n)$ 是环面 $T\subset\mathbb P_\Delta$ 的坐标。对 $\ell=(\ell_i)\in\Lambda$ 以及一个指标子集 $I$，令 $\ell_I=\sum_{i\in I}\ell_i$，计算得
\begin{equation}\label{eq:gkz-form-derivatives}
\left(\prod_{i\in I}\partial_i^{\ell_i}\right)\omega
=\frac{(-1)^{\ell_I}(\ell_I)!\prod_{i\in I}t^{\ell_im_i}}{f^{\ell_I+1}}\eta.
\end{equation}
这里在求导的指标集 $I$ 上取 $\ell_i\geq0$。由于 $\ell\in\Lambda$ 意味着
$\sum_{\ell_i>0}\ell_im_i=\sum_{\ell_i<0}(-\ell_i)m_i$，故有
\[
\prod_{\ell_i>0}t^{\ell_im_i}=\prod_{\ell_i<0}t^{-\ell_im_i}.
\]
\SourcePage{115}
再由式~\eqref{eq:gkz-form-derivatives} 和
$\sum_{\ell_i>0}\ell_i=\sum_{\ell_i<0}(-\ell_i)$，容易推出 $\Box_\ell\omega=0$。若 $\Phi$ 是 $\omega$ 的任意周期，立即得到 $\Box_\ell\Phi=0$。

同样容易证明，对 $1\leq j\leq n+1$ 有 $Z_j\Phi=0$。考虑 $(n-1)$ 形式
\[
\eta_j=\frac{dt_1}{t_1}\wedge\cdots\wedge
\widehat{\frac{dt_j}{t_j}}\wedge\cdots\wedge\frac{dt_n}{t_n},
\qquad 1\leq j\leq n.
\]
计算得
\[
d\left(\frac1f\eta_j\right)
=\frac{(-1)^j}{f^2}\left(\sum_i\lambda_im_{ij}t^{m_i}\right)\eta,
\]
其中 $m_{ij}$ 是 $m_i$ 的坐标。因为对 $1\leq j\leq n$ 有 $\widehat\beta_j=0$，容易得出 $Z_j\omega$ 对 $j<n+1$ 是恰当形式，因而这些 $j$ 都满足 $Z_j\Phi=0$。至于 $j=n+1$，有 $\widehat\beta_{n+1}=-1$，而 $(m_i,1)\in\mathcal A$ 的最后一个坐标总为 $1$，所以
\[
Z_{n+1}\omega=\left(\sum_i\lambda_i\partial_i\right)\omega+\omega
=-\frac{\sum_i\lambda_it^{m_i}}{f^2}\eta+\omega=0,
\]
正如所需。

式~\eqref{eq:toric-gkz-rational-form} 中的 $\omega$ 定义在 $T-Z_f$ 上，其中 $T\subset\mathbb P_\Delta$ 是环面，$Z_f=V\cap T$ 是由 $f\in L(\Delta)$ 确定的仿射超曲面。于是 $\omega$ 的留数是 $Z_f$ 上的 $(n-1)$ 形式。\SourceCite{Batyrev3} 更详细地解释了这一仿射观点。不过，我们需要的是 $V$ 上周期的 Picard--Fuchs 方程，而不仅是仿射部分 $Z_f$ 上的方程。这可以借助 $\mathbb P_\Delta$ 的齐次坐标环 $\mathbb C[x_1,\ldots,x_r]$ 实现。$x_i$ 与 $t_j$ 之间的关系为
\[
t_j=\prod_{i=1}^r x_i^{\langle e_j,\rho_i\rangle},
\]
在这个代换下，有
\[
\Omega_0=x_1\cdots x_r\eta,
\]
其中 $\Omega_0$ 由式~\eqref{eq:toric-volume-form} 定义。这些事实的证明见 \SourceCite{BC}。$\Omega_0$ 的公式容易理解：形式 $\eta$ 在 $\mathbb P_\Delta-T$ 的各分支 $D_1,\ldots,D_r$ 上有一阶极点。另一方面，对应于 $0\in\Delta$ 的 $\mathcal O_{\mathbb P_\Delta}(\sum_{j=1}^rD_j)$ 的截面 $x_1\cdots x_r$，在每个环面除子 $D_j$ 上恰有一阶零点；由
$\widehat\Omega_{\mathbb P_\Delta}^n\simeq\mathcal O_{\mathbb P_\Delta}(-\sum_{j=1}^rD_j)$，便得到 $\Omega_0=x_1\cdots x_r\eta$。因此，在上述变量代换下，有
\[
\omega=\frac{\eta}{f(t_1,\ldots,t_n)}
=\frac{x_1\cdots x_r\eta}{x_1\cdots x_r f(t_1,\ldots,t_n)}
=\frac{\Omega_0}{f(x_1,\ldots,x_r)},
\]
其中 $f(x_1,\ldots,x_r)$ 是按式~\eqref{eq:toric-graded-sections} 定义的 $f(t_1,\ldots,t_n)$ 的齐次化。因此，$\omega$ 是~\hyperref[sec:griffiths-dwork]{第 5.3 节}考虑过的形式之一，其留数是 $V$ 上的全纯 $(n-1)$ 形式。我们得出结论：$\mathcal A$ 超几何方程是 Picard--Fuchs 方程。这里还有一个技术问题，即 $\mathbb P_\Delta$ 可能不是单纯的；不过，一旦拉回到一个消解上，就没有困难。

$\mathcal A$ 超几何方程定义在以 $\lambda_i$ 为坐标的空间 $L(\Delta)$ 上。然而，我们关心的是模空间上的微分方程，而模空间是 $L(\Delta)$ 的一个商；更准确地说，这个商给出模的多项式部分。\hyperref[chap:moduli]{第~\ref*{chap:moduli} 章}将详细讨论此问题，现在只考察模的两个较明显的方面。

\SourcePage{116}
首先，环面 $T$ 在 $L(\Delta)$ 上有自然作用
\[
\nu\cdot f=\nu\cdot\sum_i\lambda_it^{m_i}
=\sum_i\nu^{m_i}\lambda_it^{m_i},\qquad \nu\in T.
\]
这个作用来自 $T$ 在自身上的自然作用。因此，对每个 $\nu\in T$，$f\in L(\Delta)$ 与 $\nu\cdot f$ 在 $\mathbb P_\Delta$ 中定义同构的超曲面，从而给出相同的模点。

第二个作用是 $\mathbb C^*$ 的作用，对 $c\in\mathbb C^*$ 定义为 $f\mapsto c\cdot f$。由于缩放给出同一个超曲面，$f$ 与 $c\cdot f$ 也给出相同的模点。合并这两个作用，便得到 $T\times\mathbb C^*$ 在 $L(\Delta)$ 上的作用，其商映射到模空间。

因为 $\omega=\eta/f$ 由 $f\in L(\Delta)$ 确定，$T\times\mathbb C^*$ 也作用于 $\omega$ 的周期积分。可以证明，方程 $Z_j\Phi=0$（$1\leq j\leq n$）是周期在 $T$ 的作用下不变的无穷小形式；同样，$Z_{n+1}\omega=0$ 来自缩放下的不变性。细节见 \SourceCite{HLY1}。

若限于 $\lambda_i\ne0$（对所有 $i$）的子集
$(\mathbb C^*)^{\Delta\cap M}\subset L(\Delta)$，那么
$(\mathbb C^*)^{\Delta\cap M}/(T\times\mathbb C^*)$ 很容易描述，因为它是环面的商。此时式~\eqref{eq:gkz-relation-lattice} 中的格 $\Lambda$ 就发挥作用。选取 $\Lambda$ 的一组基 $\ell_j=(\ell_{ji})$，定义函数
\[
z_j=\prod_i\lambda_i^{\ell_{ji}}.
\]
容易验证，$z_j$ 在 $T\times\mathbb C^*$ 的作用下不变，且给出商空间的坐标。我们的目标之一，是找到 $\Lambda$ 的一组基，使 $z_j$ 成为所考察的最大幺幂边界点附近复模空间的局部坐标。这样，在最大幺幂边界点附近解析的周期，就会是 $z_j$ 的解析函数。\hyperref[chap:moduli]{第~\ref*{chap:moduli} 章}将再回到这个问题。

我们希望用这些变量 $z_j$ 写出 $\mathcal A$ 超几何方程。先将方程改写成更便于处理模的形式。问题在于，$\omega$ 在 $\mathbb C^*$ 作用 $\lambda_i\mapsto c\lambda_i$ 下并不保持不变。因此，选取形式
\begin{equation}\label{eq:normalized-toric-period-form}
\widetilde\omega=\lambda_0\omega=\frac{\lambda_0}{f}\eta,
\end{equation}
它关于 $\lambda_i$ 的齐次次数现在是 $0$。因为 $\omega$ 的周期满足 $\mathcal A$ 方程组，所以 $\widetilde\omega$ 的周期 $\widetilde\Phi$ 满足
\[
Z_j\lambda_0^{-1}\widetilde\Phi=0,\qquad
\Box_\ell\lambda_0^{-1}\widetilde\Phi=0,
\]
其中 $1\leq j\leq n+1$、$\ell\in\Lambda$。

方程 $Z_j\lambda_0^{-1}\widetilde\Phi=0$ 是 $\widetilde\Phi$ 在 $T\times\mathbb C^*$ 作用下不变的无穷小形式。因此，一旦按 $T\times\mathbb C^*$ 的作用取商，余下的任务就是将方程
$\Box_\ell\lambda_0^{-1}\widetilde\Phi=0$ 改写成 $z_j$ 坐标。为此，反复使用微分算子恒等式
\begin{equation}\label{eq:logarithmic-operator-commutation}
\delta_i\lambda_i^p=\lambda_i^p(\delta_i+p).
\end{equation}
以五次三维簇的镜像说明这个方法。这里，$\Lambda$ 的相关生成元为 $\ell=(1,1,1,1,1,-5)$；回顾 $\Lambda$ 由
$\mathcal A=(\Delta\cap M)\times\{1\}$ 的元素之间的关系组成。它给出坐标
\begin{equation}\label{eq:quintic-gkz-coordinate}
z=\lambda_1\cdots\lambda_5/\lambda_0^5.
\end{equation}
\SourcePage{117}
在这种情形下，对式~\eqref{eq:toric-gkz-laurent} 作齐次化，得到
\[
\lambda_1x_1^5+\lambda_2x_2^5+\lambda_3x_3^5+\lambda_4x_4^5+\lambda_5x_5^5
+\lambda_0x_1x_2x_3x_4x_5=0.
\]
置 $\lambda_1=\cdots=\lambda_5=1$、$\lambda_0=\psi$，就得到式~\eqref{eq:quintic-pencil}。此时上面的模坐标为 $z=\psi^{-5}$，与\hyperref[chap:quintic]{第~\ref*{chap:quintic} 章}选取的坐标一致；回顾，我们现在用 $z$ 而不是 $x$ 表示模坐标。必须把 $\mathcal A$ 超几何方程
\[
\Box_\ell\lambda_0^{-1}=(\partial_1\cdots\partial_5-\partial_0^5)\lambda_0^{-1}=0
\]
改写成关于 $z$ 的方程。在 $\Box_\ell\lambda_0^{-1}$ 的左侧乘以
$\lambda_0\cdots\lambda_5=z\lambda_0^6$，得到
\[
\delta_1\cdots\delta_5-z\lambda_0^6\partial_0^5\lambda_0^{-1}
=\delta_1\cdots\delta_5-z\lambda_0^5\delta_0(\lambda_0^{-1}\delta_0)^4\lambda_0^{-1}.
\]
从右到左反复应用式~\eqref{eq:logarithmic-operator-commutation}，得到方程
\[
\delta_1\cdots\delta_5-z(\delta_0-5)(\delta_0-4)\cdots(\delta_0-1)=0.
\]
现在用 $\delta$ 表示 $z\,d/dz$。限于仅依赖 $z$ 的函数时，由式~\eqref{eq:quintic-gkz-coordinate} 得恒等式
\[
\delta_1=\cdots=\delta_5=\delta,\qquad\delta_0=-5\delta.
\]
最终得到
\begin{equation}\label{eq:quintic-gkz-fifth}
\delta^5+z(5\delta+5)(5\delta+4)\cdots(5\delta+1)=0.
\end{equation}
这种方法的局限立刻显现出来：五阶算子~\eqref{eq:quintic-gkz-fifth} 并不生成 Picard--Fuchs 理想，因为还缺少四阶算子~\eqref{eq:quintic-mum-pf}。

一般而言，从 $\mathcal A$ 方程组得到的 $z_j$ 坐标中的微分方程组，其解包含周期，但通常还会有其他解。问题在于，找到一个更大的微分方程组，其共同解恰好由周期组成。为此，有几种策略。

求出全部 Picard--Fuchs 方程的一种方法，是从 $\mathcal A$ 方程组的超几何方程出发，再添加其他 Picard--Fuchs 方程。通常，所添加的是用 Griffiths--Dwork 方法找到的低阶方程。一些例子见 \SourceCite{BKK}。

另一种常用方法，是分解来自 $\mathcal A$ 方程组的某个 Picard--Fuchs 方程。如~\hyperref[subsec:picard-fuchs]{第 5.1.2 小节}所述，只要分解得到的低阶算子 $D$ 零化所有周期，它就是 Picard--Fuchs 算子。幸而在许多情形中，只需检验 $D$ 零化一个周期 $y$ 就够了。例如，假设 $D\cdot y=0$，且把 $y$ 沿整个模空间解析延拓，会得到 Picard--Fuchs 方程解的一组基 $y_i$。那么 $D\cdot y=0$ 蕴含对所有 $i$ 都有 $D\cdot y_i=0$。因此，$D$ 零化所有周期，因而是 Picard--Fuchs 算子。

以五次三维簇的镜像说明这一过程。从五阶方程~\eqref{eq:quintic-gkz-fifth} 出发，把 $z$ 交换到右边，得到
\[
\delta^5+(5\delta)(5\delta-1)\cdots(5\delta-4)z=0.
\]
于是从左侧提出因子 $\delta$，得到四阶方程
\begin{equation}\label{eq:quintic-gkz-factor}
\delta^4+5(5\delta-1)(5\delta-2)(5\delta-3)(5\delta-4)z=0,
\end{equation}
\SourcePage{118}
它等价于式~\eqref{eq:quintic-mum-pf}。\SourceCite{HKTY1} 描述了这种分解的系统方法。为了证明它是一个 Picard--Fuchs 方程，考虑式~\eqref{eq:y0-hypergeometric} 中的
\[
y_0(z)=\sum_{n=0}^\infty\frac{(5n)!}{(n!)^5}(-1)^nz^n.
\]
由后文的式~\eqref{eq:toric-euler-period}，将看出 $y_0$ 确实是一个周期；容易检验，它是式~\eqref{eq:quintic-gkz-factor} 的解。进一步，\SourceCite{CdGP} 证明，$y_0$ 的解析延拓给出周期的一组基。因此，式~\eqref{eq:quintic-gkz-factor} 是五次三维簇镜像的 Picard--Fuchs 方程。

要在一般情形中实施这个策略，就需要知道一个满足 $\mathcal A$ 方程组的显式周期。\hyperref[chap:moduli]{第~\ref*{chap:moduli} 章}的式~\eqref{eq:toric-euler-period} 将给出这样的一个周期，取自 \SourceCite[命题 14.6]{Batyrev3}。实际上，式~\eqref{eq:toric-euler-period} 中的周期 $y$ 还给出一种略有不同的生成 Picard--Fuchs 方程的方法：利用 $y$ 的系数之间的递推关系，推出微分方程。\SourceCite{BvS} 将这个技巧用于若干单参数族。其不足是，一般尚不知道式~\eqref{eq:toric-euler-period} 的解析延拓是否给出解的一组基，尽管对 Calabi--Yau 三维簇，人们预期确实如此。更多细节见 \SourceCite{BvS}。

遗憾的是，因式分解可能仍不能给出全部缺失的 Picard--Fuchs 算子，因为 $\mathcal A$ 方程组还有第二个不完整性的来源，即自同构。GKZ 方程 $Z_j\Phi=0$ 反映的是在 $\mathbb P_\Delta$ 的环面作用下的不变性；而要得到 $f=0$ 的模，必须取商于 $\operatorname{Aut}(\mathbb P_\Delta)$ 的作用。我们在~\hyperref[sec:toric-automorphisms]{第 3.6 节}看到，$\operatorname{Aut}(\mathbb P_\Delta)$ 还包含来自根的自同构。这些自同构使某些坐标 $\lambda_i$ 无效，即对应于 $\Delta$ 的刻面内部点的那些坐标。与其添加更多方程，可以直接略去这些 $\lambda_i$；但须记住，模空间上可能还保留离散的等同关系，一个例子见 \SourceCite{CFKM}。\hyperref[chap:moduli]{第~\ref*{chap:moduli} 章}定义环面超曲面 $V$ 的\emph{简化模空间} $\mathcal M_{\mathrm{simp}}$ 时，我们将采用这一方法。另一种密切相关的方法，是把这些额外自同构的无穷小形式加入 $\mathcal A$ 方程组，见 \SourceCite{HKTY1,HLY1}。它们补充了方程 $Z_j\Phi=0$；如刚才所述，后者是环面所诱导自同构的无穷小形式。

上述 GKZ 方程组的解~\eqref{eq:toric-euler-period}，只是众多精妙的\IndexLocation{indexloc:0615:118}超几何方程解公式之一。这些公式推广了经典的超几何幂级数。\SourceCite{GKZ1,HLY1} 讨论了这一点；在\hyperref[sec:mirror-map]{第~\ref*{sec:mirror-map} 节}讨论镜像映射时，我们还将给出一些显式例子与求解方法。

再回到已经考虑过的一个例子，说明上述讨论。
\SourceExampleNumber{5.5.2}{1}
\begin{example}\label{ex:weighted-gkz}
在~\hyperref[sec:hodge-examples]{第 5.4 节}，讨论了 $\mathbb P(1,1,2,2,2)$ 中次数为 $8$ 的\IndexLocation{indexloc:1020:118}\IndexLocation{indexloc:1027:118}超曲面的镜像 $V^\circ$。这里，将式~\eqref{eq:toric-gkz-laurent} 齐次化，得到形如
\[
\lambda_1x_1^8+\lambda_2x_2^8+\lambda_3x_3^4+\lambda_4x_4^4+\lambda_5x_5^4
+\lambda_6x_1^4x_2^4+\lambda_0x_1x_2x_3x_4x_5=0
\]
的方程。这些单项式已在例题~\ref{ex:weighted-pf-method} 中算出。利用由式~\eqref{eq:weighted-pf-rays} 确定的自反多胞形 $\Delta^\circ$，可以算出 $\Lambda$ 的一组基为
$\ell_1=(0,0,1,1,1,1,-4)$ 与 $\ell_2=(1,1,0,0,0,-2,0)$。于是得到模坐标
\[
z_1=\frac{\lambda_3\lambda_4\lambda_5\lambda_6}{\lambda_0^4},\qquad
z_2=\frac{\lambda_1\lambda_2}{\lambda_6^2}.
\]
\SourcePage{119}
注意，令 $\lambda_0=\lambda_2=\lambda_4=\lambda_5=\lambda_6=1$，就有 $z_1=\lambda_3$、$z_2=\lambda_1$，与式~\eqref{eq:weighted-pf-family} 一致。由
\[
\lambda_0\lambda_1\lambda_2\Box_{\ell_2}\lambda_0^{-1}
=\lambda_0\lambda_1\lambda_2(\partial_1\partial_2-\partial_6^2)\lambda_0^{-1}=0,
\]
按上面讨论的方法，得到
\begin{equation}\label{eq:weighted-gkz-second}
\delta_{z_2}^2-z_2(\delta_{z_1}-2\delta_{z_2})(\delta_{z_1}-2\delta_{z_2}-1)=0,
\end{equation}
其中 $\delta_{z_1}=z_1\partial/\partial z_1$、$\delta_{z_2}=z_2\partial/\partial z_2$，并使用了在 $z_1,z_2$ 的函数上成立的恒等式
\begin{equation}\label{eq:weighted-logarithmic-identities}
\delta_1=\delta_2=\delta_{z_2},\qquad\delta_6=\delta_{z_1}-2\delta_{z_2}.
\end{equation}
读者应注意，简单的代数运算就能把式~\eqref{eq:weighted-gkz-second} 变为先前用 Griffiths--Dwork 方法求出的 Picard--Fuchs 方程~\eqref{eq:weighted-pf-second}；其中出现的 $\delta_i$，就是这里的 $\delta_{z_i}$。

类似地，利用式~\eqref{eq:weighted-logarithmic-identities}，以及
$\delta_3=\delta_4=\delta_5=\delta_{z_1}$、$\delta_0=-4\delta_{z_1}$，可以算出 GKZ 方程
\[
\lambda_0\lambda_3\lambda_4\lambda_5\lambda_6\Box_{\ell_1}\lambda_0^{-1}
=\lambda_0\lambda_3\lambda_4\lambda_5\lambda_6
(\partial_3\partial_4\partial_5\partial_6-\partial_0^4)\lambda_0^{-1}=0
\]
给出方程
\[
\delta_{z_1}^3(\delta_{z_1}-2\delta_{z_2})
-(4\delta_{z_1})(4\delta_{z_1}-1)(4\delta_{z_1}-2)(4\delta_{z_1}-3)z_1=0.
\]
提出左侧因子，便得到三阶方程
\[
\delta_{z_1}^2(\delta_{z_1}-2\delta_{z_2})
-4(4\delta_{z_1}-1)(4\delta_{z_1}-2)(4\delta_{z_1}-3)z_1=0,
\]
它等价于
\begin{equation}\label{eq:weighted-gkz-third}
\delta_{z_1}^2(\delta_{z_1}-2\delta_{z_2})
-4z_1(4\delta_{z_1}+1)(4\delta_{z_1}+2)(4\delta_{z_1}+3)=0.
\end{equation}
解析周期~\eqref{eq:toric-period-local-series} 满足这个方程；而且，如 \SourceCite{CdFKM} 所示，这个周期的解析延拓给出解的一组基。由本例之前的讨论，式~\eqref{eq:weighted-gkz-third} 确实是一个 Picard--Fuchs 方程。

两个 Picard--Fuchs 算子~\eqref{eq:weighted-gkz-second} 与~\eqref{eq:weighted-gkz-third} 的阶数分别为 $2$ 与 $3$，所确定方程组的解空间维数至多为 $6$。因为 $h^3(V^\circ)=6$，我们已经找到生成全部 Picard--Fuchs 方程的算子。
\end{example}

\subsection{nef 完全交}\label{subsec:nef-gkz}\index{nef完全交@nef 完全交|hyperpage}

至今已讨论 Picard--Fuchs 方程与 GKZ 方程组，但还没有多谈它们的解。这要等到\hyperref[sec:mirror-map]{第~\ref*{sec:mirror-map} 节}：在那里，将利用 Picard--Fuchs 方程的解构造镜像映射。不过，我们觉得值得先在这里介绍一个由\IndexLocation{indexloc:1349:119}环面完全交产生的有趣的 GKZ 方程组的解。

设 $X=\mathbb P_\Delta$ 是对应于\IndexLocation{indexloc:1218:119}自反多胞形 $\Delta\subset M_{\mathbb R}$ 的 $n$ 维环面簇，法扇为 $\Sigma$。为方便叙述，假定 $X$ 光滑；这个假设实际上可以适当减弱。我们考虑~\hyperref[subsec:nef-partitions]{第 4.3.1 小节}所研究完全交的一种推广。
\begin{definition}\label{def:nef-complete-intersection}
$X$ 上的一个 \emph{nef 完全交}，是完全交 $Y=\bigcap_{i=1}^kY_i$，其中每个 $Y_i$ 都是由整体截面生成的线丛 $\mathcal L_i$ 的一个一般截面的零点集，并且 $-(K_X+\sum_{i=1}^k\mathcal L_i)$ 为 nef。
\end{definition}

\SourcePage{120}
注意，最后一个条件略强于 $-K_Y$ 为 nef。我们要求对所有曲线 $C\subset X$ 都有
\[
\left(K_X+\sum_{i=1}^k\mathcal L_i\right)\cdot C\leq0.
\]
根据伴随公式，$-K_Y$ 为 nef 只要求这个条件对 $C\subset Y$ 中的曲线成立。$Y$ 与 $X$ 的有效曲线锥可以不同，我们将在\hyperref[subsec:enlarged-kahler-moduli]{第~\ref*{subsec:enlarged-kahler-moduli} 小节}讨论这一点。

注意，~\hyperref[subsec:nef-partitions]{第 4.3.1 小节}定义的 nef 分拆
$-K_X=\sum_{i=1}^kE_i$，自动定义一个 nef 完全交。这时 $\mathcal L_i=\mathcal O(E_i)$，所以
$-(K_X+\sum_{i=1}^k\mathcal L_i)$ 是平凡的，显然为 nef。

现在置 $\overline N=N\oplus\mathbb Z^k$，并对 $1\leq i\leq k$ 选取 $a_{i\rho}$，使
\[
\mathcal L_i=\mathcal O\left(\sum_{\rho\in\Sigma(1)}a_{i\rho}D_\rho\right).
\]
考虑点集
\begin{equation}\label{eq:nef-gkz-configuration}
\mathcal A=\{(v_\rho,a_{1\rho},\ldots,a_{k\rho}):\rho\in\Sigma(1)\}
\cup\{(0,e_1),\ldots,(0,e_k)\}\subset\overline N,
\end{equation}
其中 $e_i$ 是 $\mathbb Z^k$ 的第 $i$ 个标准基向量。注意，这推广了式~\eqref{eq:toric-gkz-configuration}；后者对应于 $k=1$、$\mathcal L_1=\mathcal O(\sum_\rho D_\rho)$ 的情形。

如往常一样，令 $r=|\Sigma(1)|$。对每个 $\rho\in\Sigma(1)$，置
$w_\rho=(v_\rho,a_{1\rho},\ldots,a_{k\rho})$；对 $1\leq i\leq k$，置 $u_i=(0,e_i)$。于是式~\eqref{eq:nef-gkz-configuration} 可写成
$\mathcal A=\{w_\rho,\rho\in\Sigma(1);\ u_i,1\leq i\leq k\}$。

点集 $\mathcal A$ 与光滑环面簇 $X$ 有如下关系。
\begin{lemma}\label{lem:nef-relation-lattice}
设 $\Lambda$ 是 $\mathcal A$ 的元素之间的关系格。则：
\begin{enumerate}[label=(\roman*)]
\item 对每个 $\beta\in H_2(X,\mathbb Z)$，在 $\overline N$ 中有关系
\[
\sum_\rho(D_\rho\cdot\beta)w_\rho-
\sum_i(c_1(\mathcal L_i)\cdot\beta)u_i=0.
\]
因此 $\beta$ 给出格 $\Lambda$ 的一个元素。
\item 所诱导的映射 $H_2(X,\mathbb Z)\to\Lambda$ 是同构。
\end{enumerate}
\end{lemma}
\begin{proof}
先注意，由正合列~\eqref{eq:toric-chow-sequence}，对所有 $m\in M$，在 $H^2(X,\mathbb Z)$ 中有
$\sum_\rho\langle m,v_\rho\rangle[D_\rho]=0$。因此，当 $\beta\in H_2(X,\mathbb Z)$ 时，
$\sum_\rho\langle m,v_\rho\rangle(D_\rho\cdot\beta)=0$。这对所有 $m$ 成立，故必有
$\sum_\rho(D_\rho\cdot\beta)v_\rho=0$（在 $N$ 中）。再结合
$c_1(\mathcal L_i)=\sum_\rho a_{i\rho}D_\rho$，以及 $w_\rho,u_i$ 的定义，容易得到所需关系。

为证明第二部分，考虑交换图
\[
\begin{array}{ccccccc}
0&\longrightarrow&\Lambda&\longrightarrow&\mathbb Z^{\mathcal A}&\longrightarrow&\overline N\\
&&\downarrow&&\downarrow&&\downarrow\\
0&\longrightarrow&H_2(X,\mathbb Z)&\longrightarrow&\mathbb Z^{\Sigma(1)}&\longrightarrow&N.
\end{array}
\]
上行按 $\Lambda$ 的定义正合，下行则是对正合列~\eqref{eq:toric-chow-sequence} 取对偶所得，因而也正合。中间的竖直映射将 $w_\rho$ 送到 $v_\rho$，将 $u_i$ 送到 $0$；最右侧的竖直映射是自然投影。最左侧的竖直映射由交换图确定。若 $N'$ 是 $\mathbb Z^{\Sigma(1)}\to N$ 的像，则容易看出，$\mathbb Z^{\mathcal A}\to\overline N$ 的像为 $N'\oplus\mathbb Z^k$。由此作图追踪，可以证明 $\Lambda\to H_2(X,\mathbb Z)$ 是同构。还可以检验，它的逆映射恰好就是引理陈述中描述的 $H_2(X,\mathbb Z)\to\Lambda$。
\end{proof}

为简单起见，现在假定 Kähler 锥 $K(X)$ 为单纯锥。再假定对偶 Mori 锥的本原射线生成元，是 $H_2(X,\mathbb Z)$ 的一组整基。\hyperref[subsec:enlarged-kahler-moduli]{第~\ref*{subsec:enlarged-kahler-moduli} 小节}将再讨论与这个假设有关的问题。这样可以选取一组整同调基
\[
\beta_1,\ldots,\beta_{r-n}\in H_2(X,\mathbb Z),
\]
使它们生成
\SourcePage{121}
Mori 锥 $M(X)\subset H_2(X,\mathbb R)$ 的整类 $M(X)_{\mathbb Z}$；这个锥由 $X$ 上有效曲线的同调类生成。在本小节余下部分，固定 $\Sigma(1)$ 的元素次序；不过，继续将 $\Sigma(1)$ 的生成元记为 $v_\rho$ 往往仍很方便。固定这一排序后，按式~\eqref{eq:nef-gkz-configuration} 中的次序，把变量 $\lambda_1,\ldots,\lambda_{r+k}$ 对应于 $\mathcal A$ 的各个生成元。对应于 $w_\rho\in\mathcal A$ 的变量，也常记为 $\lambda_\rho$。

置 $\overline T=\overline N\otimes\mathbb C^*$，引入 $\overline T$ 不变变量
\begin{equation}\label{eq:nef-gkz-q-coordinates}
q_j=(-1)^{\sum_{i=1}^k c_1(\mathcal L_i)\cdot\beta_j}
\frac{\prod_\rho\lambda_\rho^{D_\rho\cdot\beta_j}}
{\prod_{i=1}^k\lambda_{r+i}^{c_1(\mathcal L_i)\cdot\beta_j}},
\qquad 1\leq j\leq r-n.
\end{equation}
令 $T_1,\ldots,T_{r-n}$ 是与 $\beta_1,\ldots,\beta_{r-n}$ 对偶的 $K(X)$ 的生成元，定义
\[
q^\beta=\prod_jq_j^{T_j\cdot\beta},
\qquad \beta\in M(X)_{\mathbb Z}\subset H_2(X,\mathbb Z).
\]
\hyperref[chap:quantum-cohomology]{第~\ref*{chap:quantum-cohomology} 章}讨论量子上同调时，将更频繁地使用这个记号。还置 $t_j=\log q_j$，于是 $q_j=e^{t_j}$。

受 \SourceCite{Givental4} 启发，考虑上同调值表达式
\[
\widetilde I=e^{\sum_jt_jT_j}\sum_{\beta\in M(X)_{\mathbb Z}}q^\beta
\frac{\left(\prod_{i=1}^k\prod_{m=1}^{c_1(\mathcal L_i)\cdot\beta}
(c_1(\mathcal L_i)+m)\right)
\left(\prod_\rho\prod_{m=-\infty}^{0}(D_\rho+m)\right)}
{\prod_\rho\prod_{m=-\infty}^{D_\rho\cdot\beta}(D_\rho+m)}.
\]
指数 $e^{\sum_jt_jT_j}$ 是 $H^*(X)[t_1,\ldots,t_{r-n}]$ 中的多项式，求和则可视为
$H^*(X)[[q_1,\ldots,q_{r-n}]]$ 中的元素。注意，$c_1(\mathcal L_i)\cdot\beta\geq0$，因为 $\mathcal L_i$ 由整体截面生成；分子与分母中的因子 $D_\rho+m$ 除有限个以外全部抵消。这种略显繁琐的定义方式，是为了处理 $D_\rho\cdot\beta<0$ 的可能性。由于 $-(K_X+\sum_{i=1}^k\mathcal L_i)$ 为 nef，猜想当 $q_j$ 足够小时，$\widetilde I$ 的级数收敛。

令 $\mathcal V=\bigoplus_{i=1}^k\mathcal L_i$，其 Euler 类为
$\operatorname{Euler}(\mathcal V)=\prod_{i=1}^kc_1(\mathcal L_i)$。\hyperref[chap:mirror-theorem]{第~\ref*{chap:mirror-theorem} 章}将看到，上同调值形式函数
\begin{equation}\label{eq:nef-givental-iv}
I_{\mathcal V}=\operatorname{Euler}(\mathcal V)\widetilde I
\end{equation}
在 Givental 版本的环面完全交镜像定理中起重要作用，见式~\eqref{eq:toric-complete-intersection-i-function} 与定理~\ref{thm:givental-toric-mirror}。

在 $X=\mathbb P^4$、$k=1$、$\mathcal V=\mathcal L_1=\mathcal O_{\mathbb P^4}(5H)$ 的特殊情形，完全交就是五次三维簇。此时
\begin{equation}\label{eq:quintic-tilde-i}
\widetilde I=e^{t_1H}\sum_{d=0}^\infty e^{dt_1}
\frac{\prod_{m=1}^{5d}(5H+m)}{\prod_{m=1}^{d}(H+m)^5},
\end{equation}
而 $I_{\mathcal V}=5H\widetilde I$ 变为
\begin{equation}\label{eq:quintic-iv-hypergeometric}
I_{\mathcal V}=e^{t_1H}5H\sum_{d=0}^\infty e^{dt_1}
\frac{\prod_{m=1}^{5d}(5H+m)}{\prod_{m=1}^{d}(H+m)^5}.
\end{equation}
最后这个公式已在\hyperref[chap:quintic]{第~\ref*{chap:quintic} 章}作为式~\eqref{eq:iv-simple} 出现。

函数 $\widetilde I$ 与\IndexLocation{indexloc:0615:121}超几何方程有如下关系。令
$Q=\prod_{i=1}^k\lambda_{r+i}^{-1}$。
\begin{proposition}\label{prop:nef-i-gkz}
形式函数 $Q\widetilde I$ 满足与式~\eqref{eq:nef-gkz-configuration} 对应的 $\mathcal A$ 方程组，其中
$\widehat\beta=-\sum_{i=1}^ku_i=(0,\ldots,0,-1,\ldots,-1)$。
\end{proposition}
\begin{proof}
\SourcePage{122}
这里 $\widehat\beta$ 指的是定义式~\eqref{eq:gkz-euler-operators} 中算子 $Z_j$ 所用的向量。如~\hyperref[subsec:toric-gkz]{第 5.5.2 小节}所讨论的，对所有 $j$，$Z_j(\widetilde I)=0$ 等价于形式函数 $\widetilde I$ 的 $\overline T$ 不变性。$q_j$ 按构造在 $\overline T$ 下不变，$\lambda_i$ 在 $\widetilde I$ 中也仅通过 $q_j$ 出现，所以 $\widetilde I$ 是 $\overline T$ 不变的。乘以 $Q$ 后，其权为 $-\sum_i u_i$，从而对所有 $j$ 都有 $Z_j(Q\widetilde I)=0$。

现在考虑 $\beta\in H_2(X,\mathbb Z)\simeq\Lambda$ 所对应的算子 $\Box_\beta$。由引理~\ref{lem:nef-relation-lattice}，$\beta$ 给出关系
$\sum_\rho(D_\rho\cdot\beta)w_\rho-\sum_i(c_1(\mathcal L_i)\cdot\beta)u_i=0$。因此，$\Box_\beta$ 可写为
\begin{equation}\label{eq:nef-gkz-box}
\begin{split}
\Box_\beta={}&
\prod_{D_\rho\cdot\beta>0}\partial_\rho^{D_\rho\cdot\beta}
\prod_{c_1(\mathcal L_i)\cdot\beta<0}\partial_{r+i}^{-c_1(\mathcal L_i)\cdot\beta}\\
&-\prod_{D_\rho\cdot\beta<0}\partial_\rho^{-D_\rho\cdot\beta}
\prod_{c_1(\mathcal L_i)\cdot\beta>0}\partial_{r+i}^{c_1(\mathcal L_i)\cdot\beta},
\end{split}
\end{equation}
其中 $\partial_\rho=\partial/\partial\lambda_\rho$、$\partial_{r+i}=\partial/\partial\lambda_{r+i}$。需要证明 $\Box_\beta(Q\widetilde I)=0$，等价地，证明
$(Q^{-1}\Box_\beta Q)\widetilde I=0$。

如往常一样，置 $\delta_\rho=\lambda_\rho\partial_\rho$、
$\delta_{r+i}=\lambda_{r+i}\partial_{r+i}$，并置
$\delta_{q_j}=q_j\partial/\partial q_j=\partial/\partial t_j$。将前述方法用于式~\eqref{eq:nef-gkz-q-coordinates}，对 $\overline T$ 不变函数得到恒等式
\begin{equation}\label{eq:nef-logarithmic-identities}
\begin{aligned}
\delta_\rho&=\sum_j(D_\rho\cdot\beta_j)\delta_{q_j},\\
\delta_{r+i}&=-\sum_j(c_1(\mathcal L_i)\cdot\beta_j)\delta_{q_j}.
\end{aligned}
\end{equation}
另外，计算得
\begin{equation}\label{eq:nef-i-derivative}
\delta_{q_j}\left(e^{\sum_jt_jT_j}q^\beta\right)
=(T_j+T_j\cdot\beta)\left(e^{\sum_jt_jT_j}q^\beta\right).
\end{equation}
结合式~\eqref{eq:nef-logarithmic-identities} 与~\eqref{eq:nef-i-derivative}，得到
\begin{equation}\label{eq:nef-i-lambda-derivatives}
\begin{aligned}
\delta_\rho\left(e^{\sum_jt_jT_j}q^\beta\right)
&=(D_\rho+D_\rho\cdot\beta)\left(e^{\sum_jt_jT_j}q^\beta\right),\\
\delta_{r+i}\left(e^{\sum_jt_jT_j}q^\beta\right)
&=-(c_1(\mathcal L_i)+c_1(\mathcal L_i)\cdot\beta)
\left(e^{\sum_jt_jT_j}q^\beta\right).
\end{aligned}
\end{equation}
这里使用了除子恒等式 $D=\sum_j(D\cdot\beta_j)T_j$，分别取 $D=D_\rho$ 与 $D=c_1(\mathcal L_i)$；该恒等式成立，是因为 $T_j$ 与 $\beta_j$ 互为对偶基。

$\Box_\beta$ 用 $\partial_\rho,\partial_{r+i}$ 定义。为了改用对数微分算子，考虑
\[
\Box'_\beta=
\left(\prod_{D_\rho\cdot\beta>0}\lambda_\rho^{D_\rho\cdot\beta}
\prod_{c_1(\mathcal L_i)\cdot\beta<0}\lambda_{r+i}^{-c_1(\mathcal L_i)\cdot\beta}\right)
Q^{-1}\Box_\beta Q.
\]
记 $s_\beta=\sum_i c_1(\mathcal L_i)\cdot\beta$，则
\begin{equation}\label{eq:nef-normalized-box}
\begin{aligned}
\Box'_\beta={}&
\prod_{D_\rho\cdot\beta>0}\delta_\rho(\delta_\rho-1)\cdots
(\delta_\rho-D_\rho\cdot\beta+1)\\
&\quad\times\prod_{c_1(\mathcal L_i)\cdot\beta<0}
(\delta_{r+i}-1)(\delta_{r+i}-2)\cdots
(\delta_{r+i}+c_1(\mathcal L_i)\cdot\beta)\\
&-(-1)^{s_\beta}q^\beta
\prod_{D_\rho\cdot\beta<0}\delta_\rho(\delta_\rho-1)\cdots
(\delta_\rho+D_\rho\cdot\beta+1)\\
&\quad\times\prod_{c_1(\mathcal L_i)\cdot\beta>0}
(\delta_{r+i}-1)(\delta_{r+i}-2)\cdots
(\delta_{r+i}-c_1(\mathcal L_i)\cdot\beta).
\end{aligned}
\end{equation}
为写出其作用后的表达式，置
$a_\rho=D_\rho\cdot\beta$、$b_i=c_1(\mathcal L_i)\cdot\beta$，以及
$\varepsilon_\beta=(-1)^{-\sum_{b_i<0}b_i}$。利用式~\eqref{eq:nef-normalized-box} 与~\eqref{eq:nef-i-lambda-derivatives}，计算得
\begin{equation}\label{eq:nef-gkz-cancellation}
\begin{split}
\Box'_\beta\widetilde I
={}&\varepsilon_\beta e^{\sum_jt_jT_j}\Biggl(\\
&\sum_{\beta'\in M(X)_{\mathbb Z}}q^{\beta'}
\frac{\left(\prod_i\prod_{m=1}^{c_1(\mathcal L_i)\cdot\beta'-\min(b_i,0)}
(c_1(\mathcal L_i)+m)\right)
\left(\prod_\rho\prod_{m=-\infty}^0(D_\rho+m)\right)}
{\prod_\rho\prod_{m=-\infty}^{D_\rho\cdot\beta'-\max(a_\rho,0)}(D_\rho+m)}\\
-&\sum_{\beta''\in M(X)_{\mathbb Z}}q^{\beta''+\beta}
\frac{\left(\prod_i\prod_{m=1}^{c_1(\mathcal L_i)\cdot\beta''+\max(b_i,0)}
(c_1(\mathcal L_i)+m)\right)
\left(\prod_\rho\prod_{m=-\infty}^0(D_\rho+m)\right)}
{\prod_\rho\prod_{m=-\infty}^{D_\rho\cdot\beta''+\min(a_\rho,0)}(D_\rho+m)}
\Biggr).
\end{split}
\end{equation}

\SourcePage{123}
在 $\Box'_\beta\widetilde I$ 的这个表达式中，若 $\beta'=\beta''+\beta$，第一个和中对应于 $\beta'\in M(X)_{\mathbb Z}$ 的项，就与第二个和中对应于 $\beta''\in M(X)_{\mathbb Z}$ 的项抵消。抵消后剩下两类项：第一类来自第一个和，其中 $\beta'\in M(X)_{\mathbb Z}$ 而 $\beta'-\beta\notin M(X)_{\mathbb Z}$；第二类来自第二个和，其中 $\beta''\in M(X)_{\mathbb Z}$ 而 $\beta''+\beta\notin M(X)_{\mathbb Z}$。

先考察第一类。设 $\beta'\in M(X)_{\mathbb Z}$，且 $\beta'-\beta\notin M(X)_{\mathbb Z}$。那么 $e^{\sum_jt_jT_j}q^{\beta'}$ 出现在式~\eqref{eq:nef-gkz-cancellation} 的第一个和中，而不出现在第二个和中。目标是证明它的系数为零。假设 $\beta'-\beta\notin M(X)_{\mathbb Z}$，意味着存在一个丰沛除子 $D$，使 $D\cdot(\beta'-\beta)<0$。不妨假定 $D$ 是 $T$ 不变的，并写成 $D=\sum_\rho a_\rho D_\rho$，其中 $a_\rho\geq0$。令
\[
S=\{\rho\in\Sigma(1):D_\rho\cdot(\beta'-\beta)<0\}.
\]
由于 $D\cdot(\beta'-\beta)<0$，有 $S\ne\varnothing$。

我们断言，在 $H^*(X)$ 中有 $\prod_{\rho\in S}D_\rho=0$。为此，只需证明 $\bigcap_{\rho\in S}D_\rho=\varnothing$。假定该交集非空；那么它至少包含一个环面作用的不动点。若 $\sigma$ 是 $\Sigma$ 中对应的 $n$ 维锥，则对所有 $\rho\in S$ 都有 $\rho\subset\sigma$。$D=\sum a_\rho D_\rho$ 是 Cartier 除子，所以存在 $m_\sigma\in M$，使对 $\sigma$ 的所有射线 $\rho$，都有 $\langle m_\sigma,v_\rho\rangle=-a_\rho$；特别地，这对 $S$ 中的所有元素成立。这里 $v_\rho\in N$ 是射线 $\rho$ 的生成元。又因为 $D$ 丰沛，~\hyperref[sec:cones-fans]{第 3.1 节}末的丰沛性判据给出：对所有 $\rho\in\Sigma(1)$，都有 $\langle m_\sigma,v_\rho\rangle\geq-a_\rho$。于是 $D$ 与除子 $D'=D+\operatorname{div}(\chi^{m_\sigma})$ 线性等价。将 $D$ 换成 $D'$ 后，可以假定
\begin{equation}\label{eq:nef-gkz-divisor-contradiction}
D=\sum_{\rho\notin S}a_\rho D_\rho,\qquad a_\rho\geq0.
\end{equation}
两边与 $\beta'-\beta$ 相交，左边为负，右边非负，产生矛盾。断言得证。

现在考察式~\eqref{eq:nef-gkz-cancellation} 第一个和中 $e^{\sum_jt_jT_j}q^{\beta'}$ 的系数，并固定 $\rho\in S$。由定义，$D_\rho\cdot(\beta'-\beta)<0$；于是，从式~\eqref{eq:nef-gkz-cancellation} 可见，$D_\rho$ 是该系数的一个因子。这对所有 $\rho\in S$ 成立，所以系数是 $\prod_{\rho\in S}D_\rho=0$ 的倍数，因而为零。

剩下第二类，即 $\beta''\in M(X)_{\mathbb Z}$ 且 $\beta''+\beta\notin M(X)_{\mathbb Z}$ 的项。论证完全类似，予以省略。
\end{proof}

命题~\ref{prop:nef-i-gkz} 的一个直接推论，是形式函数 $I_{\mathcal V}=\operatorname{Euler}(\mathcal V)\widetilde I$（由式~\eqref{eq:nef-givental-iv} 定义）也是 $\mathcal A$ 方程组的解。

命题~\ref{prop:nef-i-gkz} 为式~\eqref{eq:nef-gkz-configuration} 中的 $\mathcal A$ 提供了有趣的 GKZ 方程组的解。也就是说，把 $\widetilde I$ 在 $H^*(X)$ 的任意一组基下展开，系数都满足 GKZ 方程。那么，它与 Picard--Fuchs 方程有什么联系？对于与 nef 分拆对应的 Calabi--Yau 完全交 $V\subset X$，~\hyperref[subsec:nef-partitions]{第 4.3.1 小节}构造了其镜像族 $V^\circ$；可以证明，$\mathcal A$ 方程组中的微分方程包含在镜像族 $V^\circ$ 的 Picard--Fuchs 方程中。这推广了~\hyperref[subsec:toric-gkz]{第 5.5.2 小节}的超曲面情形。（读者可能会担心，超曲面情形使用了 $\widehat\beta=(0,\ldots,0,-1)$，而完全交的归一化形式函数 $\widetilde I$ 的 Euler 权为 $\widehat\beta=(0,\ldots,0)$。这一差别并不重要，来自式~\eqref{eq:normalized-toric-period-form} 中全纯形式 $\widetilde\omega$ 的选择。）

\SourcePage{124}
举例说明：取 $X=\mathbb P^4$、$\mathcal V=\mathcal L_1=\mathcal O_{\mathbb P^4}(5H)$。用 $H^*(\mathbb P^4)$ 的通常基展开式~\eqref{eq:quintic-tilde-i}，得到
\begin{equation}\label{eq:quintic-i-expansion}
\widetilde I=y_0+y_1H+y_2H^2+y_3H^3+y_4H^4.
\end{equation}
在变量识别 $z=-q_1$ 下，$y_0,\ldots,y_4$ 都是五阶 GKZ 方程~\eqref{eq:quintic-gkz-fifth} 的解。

不过，已经看到，GKZ 方程组只是全部 Picard--Fuchs 方程的一部分。因此，虽然 $\widetilde I$ 满足归一化后的 GKZ 方程组，却可能不满足全部 Picard--Fuchs 方程。例如，在~\hyperref[subsec:toric-gkz]{第 5.5.2 小节}，式~\eqref{eq:quintic-gkz-fifth} 分解得到四阶方程，即五次镜像的 Picard--Fuchs 方程。\hyperref[subsec:hypergeometric-mirror-map]{第~\ref*{subsec:hypergeometric-mirror-map} 小节}将证明，式~\eqref{eq:quintic-i-expansion} 中只有 $y_0,\ldots,y_3$ 是 Picard--Fuchs 方程的解；$y_4$ 是多余的。因此考虑 $I_{\mathcal V}=\operatorname{Euler}(\mathcal V)\widetilde I$，将式~\eqref{eq:quintic-i-expansion} 变为
\begin{equation}\label{eq:quintic-iv-period-expansion}
I_{\mathcal V}=5H(y_0+y_1H+y_2H^2+y_3H^3),
\end{equation}
因为 $H^5=0$。\hyperref[subsec:hypergeometric-mirror-map]{第~\ref*{subsec:hypergeometric-mirror-map} 小节}的理论将说明，$y_0,\ldots,y_3$ 是五次镜像 Picard--Fuchs 方程解的一组基。还将看到，将 Frobenius 方法用于 Picard--Fuchs 方程，会自然产生式~\eqref{eq:quintic-tilde-i} 中 $\widetilde I$ 的公式。

更一般地，当 $\mathcal L_i$ 确定的完全交是 Calabi--Yau 簇时，式~\eqref{eq:nef-givental-iv} 中的
$I_{\mathcal V}=\prod_i c_1(\mathcal L_i)\widetilde I$ 应满足镜像的 Picard--Fuchs 方程。\SourceCite{Givental4} 无证明地提出了这一断言。在显式例子中，检验 $I_{\mathcal V}$ 的系数满足 $V^\circ$ 的 Picard--Fuchs 方程，往往并不困难；例如，五次三维簇的镜像确实如此。在这些例子中，通常通过考察各种上同调类空间的维数，并使用与 Picard--Fuchs 方程组对应的 $\mathcal D$ 模的不可约性，给出论证。

\hyperref[chap:mirror-theorem]{第~\ref*{chap:mirror-theorem} 章}首次出现式~\eqref{eq:quintic-iv-hypergeometric} 与~\eqref{eq:quintic-iv-period-expansion} 时，将采用 \SourceCite{LLY} 中更一般的形式。具体地，对 $\mathcal O_{\mathbb P^4}(5)$，式~\eqref{eq:hypergeometric-euler-data-function} 化为
\[
HG[\mathcal I(\widehat P)](t_1)
=e^{-pt_1/\hbar}\sum_{d\geq0}
\frac{\prod_{m=0}^{5d}(5p-m\hbar)}
{\prod_{m=1}^{d}\prod_{k=0}^{4}(p-\lambda_k-m\hbar)}e^{dt_1}.
\]
这个公式取值于等变上同调，$p$ 是等变超平面类。令 $\hbar=-1$，再令 $\lambda_k\to0$ 以回到通常上同调，得到
\[
\lim_{\lambda\to0}HG[\mathcal I(\widehat P)](t_1)
=e^{t_1H}5H\sum_{d\geq0}
\frac{\prod_{m=1}^{5d}(5H+m)}{\prod_{m=1}^{d}(H+m)^5}e^{dt_1},
\]
这恰好是此时的 $I_{\mathcal V}$。\SourceCite{LLY} 的作者选用“$HG$”这一记号，是为了强调这些函数的超几何解释。

这里引入的某些对象，将在后续章节中起重要作用。例如，\hyperref[subsec:enlarged-kahler-moduli]{第~\ref*{subsec:enlarged-kahler-moduli} 小节}将说明式~\eqref{eq:nef-gkz-q-coordinates} 中的 $q_j$，如何给出 Calabi--Yau 超曲面 $V\subset X$ 的 Kähler 模空间上的自然坐标。如前所述，\hyperref[chap:mirror-theorem]{第~\ref*{chap:mirror-theorem} 章}讨论镜像定理时，函数 $I_{\mathcal V}$ 也将占据显著位置。特别是例题~\ref{ex:toric-fano-and-hirzebruch-quantum-ring} 将说明，命题~\ref{prop:nef-i-gkz} 蕴含环面完全交量子上同调环中的一些非平凡关系。
\index{nef完全交@nef 完全交|hyperpage}\index{GKZ方程组@GKZ 方程组|hyperpage}\index{上同调值超几何级数|hyperpage}

\SourcePage{125}
\section{Yukawa 耦合}\label{sec:yukawa-couplings}\index{Yukawa耦合@Yukawa 耦合|hyperpage}

本节集中讨论 Calabi--Yau 三维簇这一特殊情形。如\hyperref[chap:introduction]{第~\ref*{chap:introduction} 章}，\emph{Yukawa 耦合}定义为
\[
\langle\theta_1,\theta_2,\theta_3\rangle
=\int_V\Omega\wedge(\nabla_{\theta_1}\nabla_{\theta_2}\nabla_{\theta_3}\Omega),
\]
其中 $\theta_i\in H^1(V,T_V)$，见式~\eqref{eq:b-yukawa}。将 $H^1(V,T_V)$ 视为模空间的切空间，并在最大幺幂边界点 $p$ 附近选取模坐标 $z_1,\ldots,z_r$，则 $\partial/\partial z_i$ 给出切空间的一组基。不过，Gauss--Manin 联络 $\nabla$ 在 $p$ 处有正则奇点，因此使用 $\delta_i=z_i\partial/\partial z_i$ 更方便。我们因而关心
\[
\langle\delta_i,\delta_j,\delta_k\rangle
=\int_V\Omega\wedge(\nabla_{\delta_i}\nabla_{\delta_j}\nabla_{\delta_k}\Omega)
\]
的计算。

\subsection{一维情形}\label{subsec:yukawa-one-dimensional}

先考虑复模空间为一维的 Calabi--Yau 三维簇 $V$。假定在一个最大幺幂边界点附近，坐标为 $z$。对 $\delta=z\,d/dz$，可将 Picard--Fuchs 方程写成
\[
\nabla_\delta^4\Omega+f_1(z)\nabla_\delta^3\Omega
+f_2(z)\nabla_\delta^2\Omega+f_3(z)\nabla_\delta\Omega+f_4(z)\Omega=0.
\]
这个方程可以确定 Yukawa 耦合，差一个常数因子。\hyperref[chap:quintic]{第~\ref*{chap:quintic} 章}已对五次镜像做过这项计算，过程很容易推广。Griffiths 横截性给出
$\int_V\Omega\wedge\nabla_\delta^2\Omega=0$；求导两次，便得到
\[
\int_V\nabla_\delta\Omega\wedge\nabla_\delta^3\Omega
=-\frac12\int_V\Omega\wedge\nabla_\delta^4\Omega,
\]
正如式~\eqref{eq:omega-identity} 附近的计算。令 $Y=\int_V\Omega\wedge\nabla_\delta^3\Omega$ 表示 Yukawa 耦合，则
\[
\begin{aligned}
\delta Y
&=\int_V\Omega\wedge\nabla_\delta^4\Omega
+\int_V\nabla_\delta\Omega\wedge\nabla_\delta^3\Omega\\
&=\frac12\int_V\Omega\wedge\nabla_\delta^4\Omega\\
&=-\frac12f_1(z)\int_V\Omega\wedge\nabla_\delta^3\Omega.
\end{aligned}
\]
最后一个等式来自 Picard--Fuchs 方程，以及由 Griffiths 横截性得到的
$\int_V\Omega\wedge\nabla_\delta^i\Omega=0$（$i\leq2$）。因此
\[
\delta Y=-\frac12f_1(z)Y.
\]
由此可解出 $Y$，差一个常数因子，正如五次镜像的情形。

\subsection{一般情形}\label{subsec:yukawa-general}

一般情形基本上使用相同的方法，只是代数运算更繁杂。为简单起见，假定 $V$ 是复模空间为二维的 Calabi--Yau 三维簇。给定一个最大幺幂边界点，选取坐标 $z_1,z_2$，使 $z_i=0$ 定义边界除子。置 $\delta_i=z_i\partial/\partial z_i$，定义
\[
K^{ab}=\int_V\Omega\wedge\nabla_{\delta_1}^{a}\nabla_{\delta_2}^{b}\Omega,
\qquad a+b=3.
\]
\SourcePage{126}
于是 $K^{ab}$ 是下列 Yukawa 耦合之一：
\[
\langle\delta_1,\delta_1,\delta_1\rangle,\quad
\langle\delta_1,\delta_1,\delta_2\rangle,\quad
\langle\delta_1,\delta_2,\delta_2\rangle,\quad
\langle\delta_2,\delta_2,\delta_2\rangle.
\]
另外，Griffiths 横截性在这里成为
\begin{equation}\label{eq:yukawa-transversality}
\int_V\Omega\wedge\nabla_{\delta_1}^{a}\nabla_{\delta_2}^{b}\Omega=0,
\qquad a+b\leq2.
\end{equation}
反复求导，可得到将
\begin{equation}\label{eq:yukawa-fourth-integrals}
W^{ab}=\int_V\Omega\wedge\nabla_{\delta_1}^{a}\nabla_{\delta_2}^{b}\Omega,
\qquad a+b=4,
\end{equation}
与 Yukawa 耦合 $K^{ab}$ 的偏导数联系起来的恒等式。再结合 Picard--Fuchs 方程，便得到一个关于 Yukawa 耦合的可解的一阶微分方程组。更多细节见 \SourceCite{HKTY1}。

为了具体说明，再回到一个熟悉的例子。
\SourceExampleNumber{5.6.2}{1}
\begin{example}\label{ex:weighted-yukawa}
$\mathbb P(1,1,2,2,2)$ 中次数为 $8$ 的\IndexLocation{indexloc:1020:126}\IndexLocation{indexloc:1030:126}\IndexLocation{indexloc:1033:126}超曲面的镜像 $V^\circ$，有二维复模空间，坐标 $z_1,z_2$ 取自例题~\ref{ex:weighted-gkz}。在那里求出了两个 Picard--Fuchs 方程~\eqref{eq:weighted-gkz-second} 与~\eqref{eq:weighted-gkz-third}，经代数运算可写成
\begin{equation}\label{eq:weighted-yukawa-pf2}
0=(1-4z_2)\delta_2^2-z_2(\delta_1^2-4\delta_1\delta_2-\delta_1+2\delta_2),
\end{equation}
\begin{equation}\label{eq:weighted-yukawa-pf3}
0=(1-256z_1)\delta_1^3-2\delta_1^2\delta_2
-8z_1(48\delta_1^2+22\delta_1+3).
\end{equation}
这些方程立即给出 $K^{ab}$ 之间的一些关系。例如，由式~\eqref{eq:weighted-yukawa-pf3} 与 Griffiths 横截性~\eqref{eq:yukawa-transversality}，容易得出
\[
0=(1-256z_1)K^{30}-2K^{21},
\]
即
\[
K^{21}=\frac12(1-256z_1)K^{30}.
\]
类似地，在式~\eqref{eq:weighted-yukawa-pf2} 的左侧分别乘以 $\delta_1,\delta_2$，再用 $K^{30}$ 解出其余耦合，得到
\[
\begin{aligned}
K^{12}&=\frac{z_2(-1+512z_1)}{1-4z_2}K^{30},\\
K^{03}&=\frac{z_2(1-256z_1+4z_2-3072z_1z_2)}{2(1-4z_2)^2}K^{30}.
\end{aligned}
\]
因此，只需确定一个 Yukawa 耦合 $K^{30}$。为此，求出它满足的一阶偏微分方程。首先，在式~\eqref{eq:weighted-yukawa-pf3} 的左侧乘以 $\delta_1$。按上面的方法，并利用 $\delta_1z_1=z_1(\delta_1+1)$，容易得到
\begin{equation}\label{eq:weighted-yukawa-fourth-relation}
(1-256z_1)W^{40}-2W^{31}-640z_1K^{30}=0,
\end{equation}
其中 $W^{ab}$ 见式~\eqref{eq:yukawa-fourth-integrals}。另外，反复对式~\eqref{eq:yukawa-transversality} 求导，得到恒等式
\[
\begin{aligned}
W^{40}&=2\delta_1K^{30},\\
W^{31}&=\frac32\delta_1K^{21}+\frac12\delta_2K^{30},\\
W^{22}&=\delta_1K^{12}+\delta_2K^{21}.
\end{aligned}
\]
\SourcePage{127}
\SourceCite{HKTY1} 给出了系统地写出全部这类恒等式的方法。将这些恒等式、上面的 $K^{21}$ 公式与式~\eqref{eq:weighted-yukawa-fourth-relation} 结合，得到
\[
(1-256z_1)\delta_1K^{30}-2\delta_2K^{30}-512z_1K^{30}=0.
\]
还可以在式~\eqref{eq:weighted-yukawa-pf2} 的左侧乘以 $\delta_1^2$，推出第二个方程：
\[
(1-4z_2)W^{22}-z_2(W^{40}-4W^{31}-K^{30}+2K^{21})=0.
\]
按相同的方法，得到
\[
-256z_1z_2\delta_1K^{30}
+\frac12(1-256z_1+1024z_1z_2)\delta_2K^{30}=0.
\]
（这里发生了大量抵消！）现在已有两个关于 $K^{30}$ 的一阶线性偏微分方程。经过代数运算，并利用 $\delta_i=z_i\partial/\partial z_i$，可改写为
\[
\begin{aligned}
\frac1{K^{30}}\frac{\partial K^{30}}{\partial z_1}
&=\frac{512(1-256z_1+1024z_1z_2)}{D(z_1,z_2)},\\
\frac1{K^{30}}\frac{\partial K^{30}}{\partial z_2}
&=\frac{512^2z_1^2}{D(z_1,z_2)},
\end{aligned}
\]
其中
\[
D(z_1,z_2)=(1-256z_1)^2-512^2z_1^2z_2.
\]
因此，容易得出 Yukawa 耦合 $K^{30}$ 的公式
\[
K^{30}=\frac{c}{D(z_1,z_2)},\qquad c\ne0.
\]
由此立即得到 $K^{21},K^{12},K^{03}$ 的公式。这些答案与 \SourceCite{HKTY1} 附录中的结果一致。

所有 Yukawa 耦合 $K^{ab}$ 的分母都涉及 $D(z_1,z_2)$ 与 $1-4z_2$。\hyperref[subsec:two-parameter-moduli]{第~\ref*{subsec:two-parameter-moduli} 小节}将证明，在 $V^\circ$ 的复模空间中，方程 $D(z_1,z_2)=0$ 与 $1-4z_2=0$ 定义判别式轨迹的两个分支；这里判别式轨迹由所有使定义方程~\eqref{eq:weighted-pf-family} 变为奇异的点 $(z_1,z_2)$ 组成。因此，这些多项式出现在 Yukawa 耦合的分母中，并不令人意外。
\end{example}

\subsection{归一化 Yukawa 耦合与 Gauss--Manin 联络}\label{subsec:normalized-yukawa-connection}\index{Gauss--Manin联络@Gauss--Manin 联络!归一化 Yukawa 耦合|hyperpage}

接下来的任务，是讨论归一化 Yukawa 耦合。如将看到的，这需要仔细选取 $3$ 形式 $\Omega$ 与模坐标。

先注意，$\Omega$ 远不是唯一的：可以乘以模空间上的任意处处非零函数。如\hyperref[chap:quintic]{第~\ref*{chap:quintic} 章}，采用以下归一化。在\IndexLocation{indexloc:0766:127}最大幺幂边界点，有 $\dim W_0=1$。由\IndexLocation{indexloc:0907:127}单值性权滤过 $W_\bullet$ 的定义，$W_0$ 在单值性作用下不变。取 $g_0$ 为 $W_0\cap H^3(V,\mathbb Z)$ 的生成元，则 $g_0$ 除符号外唯一确定。于是，$y_0=\langle g_0,\Omega\rangle$ 是在边界点解析的唯一周期，差一个常数因子。如~\hyperref[sec:quintic-b-model]{第 2.4 节}，将 $\Omega$ 换成 $\Omega/y_0$ 来归一化 $3$ 形式。\hyperref[chap:quintic]{第~\ref*{chap:quintic} 章}的论证表明，这会把 Yukawa 耦合 $Y$ 换成 $Y/y_0^2$，与式~\eqref{eq:yukawa-normalized} 前的讨论完全对应。与 \SourceCite{Morrison2} 一样，我们指出 $Y/y_0^2$ 与 $g_0$ 的选择无关。

归一化的下一步，是选取典范的模变量。这里将发现 Yukawa 耦合与 Gauss--Manin 联络 $\nabla$ 之间一种很有趣的关系。
\SourcePage{128}
为简单起见，先讨论复模为一维的 Calabi--Yau 三维簇，并假定处于最大幺幂边界点附近。如上，有 $3$ 形式 $\Omega$ 与周期 $y_0=\langle g_0,\Omega\rangle$。以 $\Omega/y_0$ 代替 $\Omega$ 归一化，同样将 $Y$ 变为 $Y/y_0^2$；归一化后，有 $\langle g_0,\Omega\rangle=1$。以下命题表明，只要选取适当局部坐标，这就会使 Picard--Fuchs 方程的形式发生显著变化。
\begin{proposition}\label{prop:normalized-yukawa-pf}
设 $V$ 是复模为一维的 Calabi--Yau 三维簇，假定所处的最大幺幂边界点满足~\hyperref[subsec:higher-dimensional-mum]{第 5.2.2 小节}的整性猜想。则存在局部坐标 $q$，使归一化 $3$ 形式 $\Omega$ 的 Picard--Fuchs 方程为
\[
\nabla_\delta^2\left(\frac{\nabla_\delta^2\Omega}{Y}\right)=0,
\]
其中 $\delta=2\pi iq\,d/dq$，$Y$ 是定义为
\[
Y=-\int_V\Omega\wedge\nabla_\delta^3\Omega
\]
的 Yukawa 耦合。
\end{proposition}
\begin{proof}
证明基于 \SourceCite{Deligne2} 的思想。先在最大幺幂边界点选取局部坐标 $z$。在 $z=0$ 的邻域，$\mathcal H=\mathcal F^0$ 上有 Hodge 滤过 $\mathcal F^\bullet$，而权滤过 $W_\bullet$ 给出 $\mathcal H$ 的滤过 $\mathcal W_\bullet$。在 $z=0$ 上，由式~\eqref{eq:mum-hodge-tate-splitting}，$(W_\bullet,F_{\mathrm{lim}}^\bullet)$ 是 Hodge--Tate 型的。如 \SourceCite[第 6 节]{Deligne2} 所指出的，这意味着在 $0$ 的邻域有
\[
\mathcal F^p\oplus\mathcal W_{2p-1}
=\mathcal F^p\oplus\mathcal W_{2p-2}=\mathcal H.
\]
因此，$(W_\bullet,\mathcal F^\bullet)$ 是 Hodge--Tate 型混合 Hodge 结构变分；对每个 $p$，$\mathcal F^p$ 在 $\operatorname{Gr}_{2p}^W$ 上诱导一维的 $(p,p)$ 型纯 Hodge 结构。置
$\mathcal I^{p,p}=\mathcal W_{2p}\cap\mathcal F^p$，则有自然同构
\[
\mathcal I^{p,p}\simeq\operatorname{Gr}_{2p}^W.
\]
由于 $\mathcal W_\bullet$ 平坦，$\nabla$ 在 $\operatorname{Gr}_{2p}^W$ 上诱导联络 $\nabla_{2p}^W$。又有
$N(\mathcal W_{2p})\subset\mathcal W_{2p-2}=\mathcal W_{2p-1}$，其中 $N$ 是单值性的对数，故 $N$ 在 $\operatorname{Gr}_{2p}^W$ 上作用为零。于是 $\nabla_{2p}^W$ 在 $\operatorname{Gr}_{2p}^W$ 上的单值性平凡，自然延拓到 $z=0$。$\operatorname{Gr}_{2p}^W$ 为一维，且有自然的 $\mathbb Z$ 结构，因此在 $0$ 的邻域可以找到张成它的整值 $\nabla_{2p}^W$ 平坦截面，除 $\pm1$ 外唯一。借助上述同构，这给出 $\mathcal I^{p,p}$ 的一个截面，记为 $e_{3-p}$。

截面 $e_0,e_1,e_2,e_3$ 有一些很好的性质。首先，$N$ 是配对
$\langle\alpha,\beta\rangle=\int_V\alpha\wedge\beta$ 的无穷小自同构。结合式~\eqref{eq:monodromy-weight-filtration} 与 $\mathcal W_4=\mathcal W_5$，得到
$\langle\mathcal W_0,\mathcal W_4\rangle=\langle\mathcal W_2,\mathcal W_2\rangle=\{0\}$。因此，只有
$\langle e_0,e_3\rangle$ 与 $\langle e_1,e_2\rangle$ 可以非零，并可在 $\operatorname{Gr}_\bullet^W$ 中计算。配对在 $\mathbb Z$ 上完美，故它们必为 $\pm1$。可以选取 $e_0,e_1,e_2,e_3$，使
\begin{equation}\label{eq:normalized-integral-pairings}
\langle e_0,e_3\rangle=-1,\qquad\langle e_1,e_2\rangle=1.
\end{equation}
这个符号选择的原因很快就会清楚。

\SourcePage{129}
现在令 $\delta=2\pi iz\,d/dz$。因为 $e_{3-p}$ 映到 $\operatorname{Gr}_{2p}^W$ 的一个 $\nabla_{2p}^W$ 平坦截面，$\nabla(e_{3-p})$ 属于 $\mathcal W_{2p-1}=\mathcal W_{2p-2}$。另一方面，$e_{3-p}$ 是 $\mathcal F^p$ 的截面，所以 $\nabla_\delta(e_{3-p})$ 属于 $\mathcal F^{p-1}$。因此，它是 $\mathcal I^{p-1,p-1}$ 的截面，从而对某些函数 $Y_1,Y_2,Y_3$，有
\[
\nabla_\delta e_0=Y_1e_1,\quad
\nabla_\delta e_1=Y_2e_2,\quad
\nabla_\delta e_2=Y_3e_3,\quad
\nabla_\delta e_3=0.
\]
由于 $\langle e_0,e_2\rangle$ 恒为零，
\[
\begin{aligned}
0=\delta\langle e_0,e_2\rangle
&=\langle\nabla_\delta e_0,e_2\rangle+\langle e_0,\nabla_\delta e_2\rangle\\
&=\langle Y_1e_1,e_2\rangle+\langle e_0,Y_3e_3\rangle=Y_1-Y_3,
\end{aligned}
\]
最后一步用了式~\eqref{eq:normalized-integral-pairings}，故 $Y_1=Y_3$。

$\nabla$ 有正则奇点，而 $\delta=2\pi iz\,d/dz$，所以 $Y_3$ 在 $0$ 处全纯。下面证明 $Y_3(0)=\pm1$。由式~\eqref{eq:connection-monodromy-residue}，在 $0$ 上，把它视为
$F_{\mathrm{lim}}^1/F_{\mathrm{lim}}^2\to F_{\mathrm{lim}}^0/F_{\mathrm{lim}}^1$ 的映射时，有
\[
\overline\nabla_\delta=-N.
\]
于是 $\overline\nabla_\delta(e_2)=-N(e_2)$。又由
$\nabla_\delta(e_2)=Y_3e_3$，有
$\overline\nabla_\delta(e_2)=Y_3(0)e_3$，因而
$N(e_2)=-Y_3(0)e_3$。但在 $0$ 上，$e_3$ 是 $W_0$ 的一组整基，$e_2,e_3$ 是 $W_2$ 的一组整基。由~\hyperref[subsec:higher-dimensional-mum]{第 5.2.2 小节}末的整性猜想，$N(e_2)=\pm e_3$。所以 $Y_3(0)=\pm1$，正如所需。

必要时将 $e_1,e_2$ 换成 $-e_1,-e_2$，便可以假定 $N(e_2)=-e_3$；注意，式~\eqref{eq:normalized-integral-pairings} 不变。于是 $Y_3(0)=1$，函数
\[
q=\exp\left(\int Y_3(z)\frac{dz}{z}\right)
\]
给出 $0$ 处的局部坐标，满足
\[
Y_3(z)\,2\pi iq\frac d{dq}=2\pi iz\frac d{dz}.
\]
重新将 $\delta$ 定义为 $2\pi iq\,d/dq$，则 $\nabla_\delta(e_2)=e_3$。因为 $Y_1=Y_3$，也有 $\nabla_\delta(e_0)=e_1$。因此，在基 $e_0,e_1,e_2,e_3$ 下，以 $q$ 为局部坐标，$\nabla_\delta$ 的\IndexLocation{indexloc:0174:129}联络矩阵为
\begin{equation}\label{eq:normalized-yukawa-matrix}
\begin{pmatrix}
0&0&0&0\\1&0&0&0\\0&Y&0&0\\0&0&1&0
\end{pmatrix},
\end{equation}
其中 $Y=Y_2$。最后将把 $Y$ 与归一化 Yukawa 耦合等同。

下一步证明，$e_0$ 就是归一化 $3$ 形式 $\Omega$。$e_0$ 是 $\mathcal F^3$ 的处处非零截面，因此对某个全纯函数 $a(q)$，有 $\Omega=a(q)e_0$。上述矩阵给出
\[
\nabla_\delta^j\Omega=\delta^j(a)e_0\pmod{e_1,e_2,e_3},
\]
所以 $a(q)$ 满足与 $\Omega$ 相同的 Picard--Fuchs 方程。在最大幺幂边界点，除常数倍外只有一个全纯解；由于 $\Omega$ 已归一化，$\langle g_0,\Omega\rangle=1$ 就是一个解。这说明 $a(q)$ 是常数 $a$，故 $\Omega=ae_0$。回顾，$g_0$ 是 $W_0$ 的整基，所以可以假定 $e_3=g_0$。配对为反对称，因而
\[
1=\langle g_0,\Omega\rangle=\langle e_3,ae_0\rangle
=-a\langle e_0,e_3\rangle=a,
\]
\SourcePage{130}
其中最后一个等号来自式~\eqref{eq:normalized-integral-pairings}。这就证明了 $e_0=\Omega$。由此很容易得到所需的 Picard--Fuchs 方程：
\begin{equation}\label{eq:normalized-yukawa-pf}
\nabla_\delta^2\left(\frac{\nabla_\delta^2\Omega}{Y}\right)=0.
\end{equation}
证明的最后一步是说明 $Y$ 就是 Yukawa 耦合。这很容易，因为已经证明 $\Omega=e_0$；于是式~\eqref{eq:normalized-yukawa-matrix} 给出 $\nabla_\delta^3\Omega=\delta(Y)e_2+Ye_3$，再由式~\eqref{eq:normalized-integral-pairings} 得到
\[
\int_V\Omega\wedge\nabla_\delta^3\Omega
=\int_V e_0\smile\bigl(\delta(Y)e_2+Ye_3\bigr)
=\langle e_0,\delta(Y)e_2+Ye_3\rangle=-Y.
\]
命题得证。
\end{proof}

证明过程中如此仔细地处理了所有符号，最后 $Y$ 的公式却出现一个负号，这可能令人觉得奇怪。不过，这个负号完全自然，因为在等式
\begin{equation}\label{eq:normalized-yukawa-sign}
Y=-\int_V\Omega\wedge\nabla_\delta^3\Omega
\end{equation}
中，右端就是使 $H^3(V,\mathbb C)$ 上的 Hodge 结构极化的配对
$Q(\alpha,\beta)=(-1)^{3(3-1)/2}\int_V\alpha\smile\beta$
（见式~\eqref{eq:hodge-polarization}）。按照镜像对称，$Y$ 应对应镜像 $V^\circ$ 上的某个 A 模型关联函数；后者也可以用 \emph{A 模型 Hodge 结构变分}上的极化来定义，我们将在\hyperref[sec:a-model-vhs]{第~\ref*{sec:a-model-vhs} 节}讨论这一点。比较这两个极化 Hodge 结构变分时，就会明白为什么需要这个符号。用\hyperref[subsec:normalized-yukawa-mirror]{第 5.6.4 小节}的术语来说，式~\eqref{eq:normalized-yukawa-sign} 是一维模空间情形的\emph{归一化 Yukawa 耦合}。

\SourceCite{Deligne2} 中对命题~\ref{prop:normalized-yukawa-pf} 的证明比上述证明更为精巧。Deligne 从原来的局部坐标 $z$ 出发，固定 $1\leq i\leq3$，考虑 Hodge 结构的扩张
\begin{equation}\label{eq:hodge-extension-coordinates}
0\longrightarrow\operatorname{Gr}_{2i-2}^W
\longrightarrow\mathcal W_{2i}/\mathcal W_{2i-4}
\longrightarrow\operatorname{Gr}_{2i}^W\longrightarrow0.
\end{equation}
在每一点，所得扩张由 $\mathbb C^*$ 中的一个数分类。因此，当 $(\mathcal W_\bullet,\mathcal F^\bullet)$ 在 $z=0$ 的一个邻域内变化时，扩张就由一个函数 $q_i(z)$ 分类。前面记为 $Y_1,Y_2,Y_3$ 的函数可用 $q_i$ 表示为
\[
Y_{4-i}=\frac{z}{q_i}\frac{dq_i}{dz},\qquad1\leq i\leq3.
\]
（\SourceCite{Deligne2} 的第 7 节中出现一个负号，因为那里处理的是对偶情形。）所以命题~\ref{prop:normalized-yukawa-pf} 中的特殊坐标 $q$ 就是 $q_1=q_3$，它分别分类 $i=1,3$ 时的扩张~\eqref{eq:hodge-extension-coordinates}。

命题~\ref{prop:normalized-yukawa-pf} 可以很好地推广到 Calabi--Yau 三维簇具有 $r$ 维模空间的情形。假定\hyperref[subsec:higher-dimensional-mum]{第 5.2.2 小节}的整性猜想成立，那么在一个最大幺幂边界点附近，就可以如下描述 Gauss--Manin 联络。我们把该边界点视为 $0\in\Delta^r$。

与前面一样，$I^{p,p}=\mathcal W_{2p}\cap\mathcal F^p\simeq\operatorname{Gr}_{2p}^W$，不过现在 $I^{1,1}$ 和 $I^{2,2}$ 的秩为 $r$。命题证明中使用的基 $e_0,e_1,e_2,e_3$ 相应地推广为以下各空间的基：
\[
e_0\in I^{3,3},\quad e_j\in I^{2,2}\ (1\leq j\leq r),\quad
 e^k\in I^{1,1}\ (1\leq k\leq r),\quad e^0\in I^{0,0},
\]
使得
\begin{equation}\label{eq:multimoduli-integral-pairings}
\langle e_0,e^0\rangle=-1,\qquad
\langle e_j,e^k\rangle=\delta_{jk},\quad1\leq j,k\leq r,
\end{equation}
\SourcePage{131}
其余配对全部为零。注意，式~\eqref{eq:normalized-integral-pairings} 正是 $r=1$ 的特殊情形。此外，与命题~\ref{prop:normalized-yukawa-pf} 的证明中一样，可以假定这些基元素延拓到 $0$，并且在那里具有整性。

证明命题~\ref{prop:normalized-yukawa-pf} 所用的方法（或 \SourceCite{Deligne2} 中更强的方法）于是表明，可以找到 $0\in\Delta^r$ 的特殊局部坐标 $q_1,\ldots,q_r$，使得对 $\delta_j=2\pi i q_j\partial/\partial q_j$，有
\begin{equation}\label{eq:multimoduli-connection}
\begin{aligned}
\nabla_{\delta_i}e^0&=0,\\
\nabla_{\delta_i}e^k&=\delta_{ik}e^0,&&1\leq k\leq r,\\
\nabla_{\delta_i}e_j&=\sum_{k=1}^rY_{ijk}e^k,&&1\leq j\leq r,\\
\nabla_{\delta_i}e_0&=e_i,
\end{aligned}
\end{equation}
其中 $Y_{ijk}$ 是一些函数。这说明了矩阵~\eqref{eq:normalized-yukawa-matrix} 如何推广到 $r$ 维模空间的情形。

由命题~\ref{prop:normalized-yukawa-pf} 的论证，知道 $e_0$ 就是归一化 $3$ 形式 $\Omega$，于是式~\eqref{eq:multimoduli-connection} 表明
\begin{equation}\label{eq:multimoduli-yukawa}
Y_{ijk}=-\int_V\Omega\wedge\nabla_{\delta_k}\nabla_{\delta_i}\nabla_{\delta_j}\Omega.
\end{equation}
在\hyperref[subsec:normalized-yukawa-mirror]{第 5.6.4 小节}中，将把 $Y_{ijk}$ 定义为\emph{归一化 Yukawa 耦合}。上述讨论于是说明，Gauss--Manin 联络由归一化 Yukawa 耦合唯一确定。

下面证明，函数 $Y_{ijk}$ 都来自同一个“势函数”。引入变量 $t_j$，使 $q_j=e^{2\pi i t_j}$，于是 $\partial/\partial t_j=2\pi i q_j\partial/\partial q_j=\delta_j$。有如下引理。
\begin{lemma}\label{lem:yukawa-potential}\index{Yukawa 势函数|hyperpage}
存在全纯函数 $\Phi(t_1,\ldots,t_r)$，使得对所有 $i,j,k$，
\[
Y_{ijk}=\frac{1}{(2\pi i)^3}\frac{\partial^3\Phi}{\partial t_i\partial t_j\partial t_k}.
\]
\end{lemma}
\begin{proof}
关键是 $\nabla$ 为平坦联络，因此复合映射
\[
\mathcal H\xrightarrow{\nabla}\mathcal H\otimes\Omega_S^1
\xrightarrow{\nabla}\mathcal H\otimes\Omega_S^2
\]
为零（这里 $S=(\Delta^*)^r$）。利用式~\eqref{eq:multimoduli-connection}，计算得
\[
0=\nabla^2(e_0)=\sum_{i,j,k=1}^rY_{ijk}e^k\otimes dt_i\wedge dt_j.
\]
这表明，对所有 $i,j,k$，$Y_{ijk}=Y_{jik}$。此外，对 $1\leq j\leq r$，又可计算出
\[
0=\nabla^2(e_j)=\sum_{i,k=1}^rY_{ijk}e^0\otimes dt_k\wedge dt_i
+\sum_{i,k,\ell=1}^r\frac{\partial Y_{ijk}}{\partial t_\ell}e^k\otimes dt_\ell\wedge dt_i.
\]
这首先表明，对所有 $i,j,k$，$Y_{ijk}=Y_{kji}$；其次表明，对所有 $i,j,k,\ell$，
\[
\frac{\partial Y_{ijk}}{\partial t_\ell}
=\frac{\partial Y_{\ell jk}}{\partial t_i}.
\]
所以 $Y_{ijk}$ 关于 $i,j,k$ 完全对称。再利用上述偏导数条件，很容易推出 $\Phi$ 的存在性。
\end{proof}

\SourcePage{132}
注意，如果用式~\eqref{eq:multimoduli-connection} 定义 $\nabla$，那么势函数 $\Phi$ 的存在性等价于 $\nabla$ 的平坦性。

这样就证明了：在满足整性猜想的最大幺幂边界点 $0\in\Delta^r$ 附近，Calabi--Yau 三维簇族的 Gauss--Manin 联络由一个函数 $\Phi$ 确定。由于 $\Phi$ 是 $t_j=\frac1{2\pi i}\log q_j$ 的函数，它在 $(\Delta^*)^r$ 上是多值的；不过，它的三阶导数就是归一化 Yukawa 耦合，而后者是单值的，并全纯延拓到 $\Delta^r$。此外，利用 $\Phi$，还可以显式写出 $H^3(V,\mathbb C)$ 上的 Hodge 结构变分的许多有趣性质。例如，在式~\eqref{eq:a-model-flat-sections-threefold} 中将看到，可以用势函数 $\Phi$ 写出一组多值平坦截面的基。

真正令人惊讶的一点是，镜像一侧存在 $\Phi$ 的一个自然对应物，称为 \emph{Gromov--Witten 势}。将在\hyperref[sec:a-model-vhs]{第~\ref*{sec:a-model-vhs} 节}和\hyperref[sec:mirror-symmetry-revisited]{第~\ref*{sec:mirror-symmetry-revisited} 节}详细探究这种对应关系。

\subsection{归一化 Yukawa 耦合与镜像映射}\label{subsec:normalized-yukawa-mirror}
命题~\ref{prop:normalized-yukawa-pf} 表明，只要使用归一化 $3$ 形式及正确的局部坐标，Gauss--Manin 联络就由 Yukawa 耦合确定。现在出现了一个奇迹：\emph{那个神秘的坐标 $q$ 正是镜像映射！}镜像映射的一般定义要到\hyperref[sec:mirror-map]{第~\ref*{sec:mirror-map} 节}才给出，不过，沿用\hyperref[chap:quintic]{第 2 章}的做法，就能理解一维模空间的情形。回顾该章公式~\eqref{eq:mirror-map-periods}，它将镜像映射写成
\begin{equation}\label{eq:mirror-map-period-ratio}
\widehat q=\exp\left(2\pi i\frac{\int_{\gamma_1}\Omega}{\int_{\gamma_0}\Omega}\right),
\end{equation}
其中 $\int_{\gamma_0}\Omega$ 为全纯函数，而绕 $z=0$ 一周时，$\int_{\gamma_1}\Omega/\int_{\gamma_0}\Omega$ 增加 $1$。令 $g_0,g_1$ 为 $\gamma_0,\gamma_1$ 的 Poincaré 对偶，就得到 $g_0\in W_0$、$g_1\in W_2$ 以及 $N(g_1)=g_0$，其中 $N$ 为单值性的对数。所以 $g_0,g_1$ 满足最大幺幂单值性定义中的条件，而 $\widehat q$ 的公式变成
\begin{equation}\label{eq:mirror-map-normalized-period}
\widehat q=\exp\bigl(2\pi i\langle g_1,\Omega\rangle\bigr),
\end{equation}
因为 $\Omega$ 的归一化意味着 $\langle g_0,\Omega\rangle=1$。

可以如下证明，这就是命题~\ref{prop:normalized-yukawa-pf} 中的坐标 $q$。取命题证明中构造的基 $e_0,e_1,e_2,e_3$；该证明还给出了 $g_0=e_3$、$\Omega=e_0$ 和 $\nabla_\delta(e_2)=e_3$。由此很容易看出
\[
g_1=-e_2+\frac{\log q}{2\pi i}e_3
\]
为平坦截面，这里的 $q$ 是命题~\ref{prop:normalized-yukawa-pf} 中的坐标。逆时针绕原点一周，$\frac{\log q}{2\pi i}$ 增加 $1$，所以 $\mathcal T(g_1)=g_1+g_0$，从而 $N(g_1)=g_0$。按定义，$e_2,e_3$ 分别属于 $W_2,W_0$，于是 $g_1$ 是 $W_2$ 中满足 $N(g_1)=g_0$ 的平坦截面。这就是所需的 $g_1$，由此得到
\[
\langle g_1,\Omega\rangle=\frac1{2\pi i}\log q.
\]
因此，坐标 $q$ 就是式~\eqref{eq:mirror-map-normalized-period} 定义的镜像映射 $\widehat q$。

更一般地，在\hyperref[chap:moduli]{第~\ref*{chap:moduli} 章}给出镜像映射的一般定义以后，将会看到：对具有 $r$ 维模空间的 Calabi--Yau 三维簇，镜像映射正是前面推广命题~\ref{prop:normalized-yukawa-pf} 时，在式~\eqref{eq:multimoduli-integral-pairings} 和式~\eqref{eq:multimoduli-connection} 的讨论中出现的特殊坐标 $q_1,\ldots,q_r$。

\SourcePage{133}
暂且接受这一结论，将 $q_1,\ldots,q_r$ 称为\emph{镜像坐标}；利用这些坐标，终于可以给出归一化 Yukawa 耦合的正式定义。
\begin{definition}\label{def:normalized-yukawa}
设有一个 Calabi--Yau 三维簇族，它在某个\IndexLocation{indexloc:0766:133}最大幺幂边界点满足整性猜想。若 $q_1,\ldots,q_r$ 是上述镜像坐标，$\Omega$ 是归一化 $3$ 形式，则\emph{归一化 Yukawa 耦合}定义为
\[
Y_{ijk}=-\int_V\Omega\wedge\nabla_{\delta_k}\nabla_{\delta_i}\nabla_{\delta_j}\Omega,
\qquad \delta_j=2\pi i q_j\frac{\partial}{\partial q_j}.
\]
\end{definition}
\index{Yukawa耦合@Yukawa 耦合!归一化|hyperpage}\index{镜像坐标|hyperpage}
$Y_{ijk}$ 也常称为 \emph{B 模型关联函数}。有一个情形，可以显式计算归一化 Yukawa 耦合。

\SourceExampleNumber{5.6.4}{1}
\begin{example}\label{ex:quintic-normalized-yukawa}
令 $V^\circ$ 为五次三维簇的镜像。由 \SourceCite{Morrison2}，知道 $V^\circ$ 满足整性猜想，所以它的归一化 Yukawa 耦合 $Y$ 有定义。我们断言，$Y$ 由下式给出：
\begin{equation}\label{eq:quintic-normalized-yukawa}
Y=\frac{5}{(1+5^5z)y_0(z)^2}\left(\frac qz\frac{dz}{dq}\right)^3,
\end{equation}
其中 $z$ 是通常的模变量，且
\[
\begin{aligned}
y_0(z)&=\sum_{n=0}^\infty\frac{(5n)!}{(n!)^5}(-1)^nz^n,\\
y_1(z)&=y_0(z)\log(-z)+5\sum_{n=1}^\infty\frac{(5n)!}{(n!)^5}
\left[\sum_{j=n+1}^{5n}\frac1j\right](-1)^nz^n,\\
q&=\exp(y_1/y_0).
\end{aligned}
\]
注意，式~\eqref{eq:quintic-normalized-yukawa} 正是\hyperref[chap:quintic]{第 2 章}五次三维簇镜像对称陈述~\eqref{eq:quintic-mirror-identity} 中出现的 B 模型关联函数。

\hyperref[subsec:hypergeometric-mirror-map]{第~\ref*{subsec:hypergeometric-mirror-map} 小节}将证明 $y_0,y_1$ 是\IndexLocation{indexloc:1176:133}\IndexLocation{indexloc:1178:133}\IndexLocation{indexloc:1180:133}五次镜像的 Picard--Fuchs 方程的解。\SourceCite{CdGP} 用公式 $q=\exp(y_1/y_0)$ 定义镜像映射，而 \SourceCite{Morrison2} 附录 C 证明了它与式~\eqref{eq:mirror-map-period-ratio} 的定义一致。因此 $q$ 就是镜像坐标。

若 $\Omega$ 为式~\eqref{eq:quintic-residue} 定义的 $3$ 形式，那么将\hyperref[subsec:yukawa-one-dimensional]{第 5.6.1 小节}的方法应用于\hyperref[ex:quintic-pf-method]{例题~\ref*{ex:quintic-pf-method}} 中的 Picard--Fuchs 方程~\eqref{eq:quintic-mum-pf}，便得到
\[
\int_{V^\circ}\Omega\wedge\nabla_{z\frac d{dz}}\nabla_{z\frac d{dz}}\nabla_{z\frac d{dz}}\Omega
=\frac1{1+5^5z}\times\text{某个常数},
\]
与\hyperref[chap:quintic]{第~\ref*{chap:quintic} 章}的结果相同。归一化就是用 $\widetilde\Omega=\Omega/\langle g_0,\Omega\rangle$ 替代 $\Omega$。由于 $y_0$ 和 $\langle g_0,\Omega\rangle$ 都是 Picard--Fuchs 方程的全纯解，$y_0$ 是 $\langle g_0,\Omega\rangle$ 的一个非零常数倍，因此，如\hyperref[subsec:normalized-yukawa-connection]{第~\ref*{subsec:normalized-yukawa-connection} 小节}所解释的，
\[
\int_{V^\circ}\widetilde\Omega\wedge\nabla_{z\frac d{dz}}\nabla_{z\frac d{dz}}\nabla_{z\frac d{dz}}\widetilde\Omega
=\frac1{(1+5^5z)y_0(z)^2}\times\text{另一个常数}.
\]
最后，从 $z$ 换到镜像坐标 $q$ 时，有
\[
\delta=2\pi i q\frac d{dq}=2\pi i\frac qz\frac{dz}{dq}\,z\frac d{dz},
\]
\SourcePage{134}
所以由 Griffiths 横截性，归一化 Yukawa 耦合为
\begin{equation}\label{eq:quintic-yukawa-undetermined}
\begin{aligned}
Y&=-\int_{V^\circ}\widetilde\Omega\wedge\nabla_\delta\nabla_\delta\nabla_\delta\widetilde\Omega\\
&=\frac{c}{(1+5^5z)y_0(z)^2}\left(\frac qz\frac{dz}{dq}\right)^3,
\end{aligned}
\end{equation}
其中 $c$ 是某个非零常数。因此，只要证明 $c=5$，就得到式~\eqref{eq:quintic-normalized-yukawa}。在 $q=0$ 处取上式的值，得到 $Y(0)=c$，所以只需证明 $Y(0)=5$。

用\hyperref[chap:quintic]{第~\ref*{chap:quintic} 章}的记号，常数 $c$ 被写成 $(2\pi i)^3c_2$；在式~\eqref{eq:quintic-mirror-yukawa} 后面的讨论中，我们令 $c_2=5/(2\pi i)^3$，以得到镜像对称所预言的正确首项。不过，在本章中，可以通过计算归一化 Yukawa 耦合，\emph{预先预测这个常数}。更确切地说，将结合本章前面建立的 Hodge 理论与 \SourceCite{CdGP,Morrison2} 中的显式计算，求出 $Y(0)$。

需要使用命题~\ref{prop:normalized-yukawa-pf} 证明中的记号。回顾，$e_0=\widetilde\Omega$ 是典范延拓 $\widetilde{\mathcal H}$ 的一组基 $e_0,e_1,e_2,e_3$ 的一部分，并且这组基在 $q=0$ 处是整结构的一组基。

现在考察 $q=0$ 处的 $\overline\nabla_\delta^3(e_0)$。首先，由式~\eqref{eq:connection-monodromy-residue}，知道在 $q=0$ 处，极限 Hodge 滤过的分次部分上有
\[
\overline\nabla_\delta=-N,
\]
所以模去 $F_{\mathrm{lim}}^1$ 后，$\overline\nabla_\delta^3(e_0)(0)=(-N)^3(e_0(0))$。另一方面，由联络矩阵~\eqref{eq:normalized-yukawa-matrix}，还有 $\overline\nabla_\delta^3(e_0)=Ye_3+\delta(Y)e_2$。合起来便得到
\[
N^3(e_0(0))=-Y(0)e_3(0).
\]
令 $s_i$ 为 $\mathcal H$ 的平坦截面，使得 $\exp(-t_1N)s_i$ 在 $q=0$ 处的值为 $e_i(0)$。由于 $e_i$ 在 $q=0$ 处是一组整基，$s_i$ 就是 $\mathcal H_{\mathbb Z}$ 的一组整基。由前一等式得
\begin{equation}\label{eq:quintic-cubic-monodromy}
N^3(s_0)=-Y(0)s_3.
\end{equation}
最后，使用 \SourceCite{CdGP} 中\IndexLocation{indexloc:1174:134}\IndexLocation{indexloc:1180:134}五次镜像族单值性的显式计算。如 \SourceCite{Morrison2} 附录 C 所解释的，$\mathcal H$ 有整平坦截面 $\beta^2,\alpha_2$，它们属于一组 $\mathbb Z$ 基，满足 $\langle\alpha_2,\beta^2\rangle=1$ 及 $N^3(\beta^2)=-5\alpha_2$。（\SourceCite{Morrison2} 中写成正号，因为那里的边界点是 $\infty$，这里却是 $0$；把 $\infty$ 变到 $0$ 的反演映射反转了相应单值性群生成元的定向。）

比较式~\eqref{eq:quintic-cubic-monodromy} 与 $N^3(\beta^2)=-5\alpha_2$，知道 $s_3,\alpha_2$ 是 $W_0$ 的整生成元，因此必要时同时将 $\beta^2,\alpha_2$ 乘以 $-1$，即可假定 $\alpha_2=s_3$。由于 $N^3:W_6/W_5\simeq W_0$，且 $N^3(Y(0)\beta^2-5s_0)=0$，推知 $Y(0)\beta^2-5s_0\in W_5$。于是 $\langle W_5,W_0\rangle=0$ 表明
$\langle Y(0)\beta^2-5s_0,s_3\rangle=0$。另一方面，已经知道 $\langle\alpha_2,\beta^2\rangle=1$，而式~\eqref{eq:normalized-integral-pairings} 给出 $\langle s_0,s_3\rangle=-1$。立刻得到 $Y(0)=5$，这正是所需的结论。
\end{example}

关于刚才的例题，还想补充两点。首先，$Y(0)=5$ 与式~\eqref{eq:quintic-cubic-monodromy} 表明
\[
N^3(s_0)=-5s_3.
\]
从镜像对称的角度看，这一等式完全合乎情理。在\hyperref[sec:a-model-vhs]{第~\ref*{sec:a-model-vhs} 节}研究五次三维簇的 A 模型联络时，将看到该联络的单值性 $N$ 自然同构于
\SourcePage{135}
与 $-H$ 的杯积，其中 $H$ 是五次三维簇的超平面截面。所以上述单值性等式与 $H^3=5$ 一致。

其次，证明式~\eqref{eq:quintic-normalized-yukawa} 需要对五次镜像族有\emph{极为详尽的了解}。除本章讨论的计算外，还需要明确的单值性信息，才能验证整性猜想、计算镜像映射并确定 $Y(0)$。

这番讨论说明：归一化 Yukawa 耦合的定义虽然有很好的理论性质，实际计算却十分困难。更常见的做法是，只对 Yukawa 耦合进行部分归一化，即将 $3$ 形式 $\Omega$ 归一化，并换成镜像坐标，而暂不确定在 $0$ 处的值。在这种情形中，只能将 $Y_{ijk}$ 确定到一个常数；不过，由引理~\ref{lem:yukawa-potential}，需要处理的常数只有一个。在表述镜像定理的环面版本时，将采取这一做法。

本节从 Yukawa 耦合开始，随后讨论 Picard--Fuchs 方程及联络矩阵，最后引出了镜像映射和归一化 Yukawa 耦合的概念。这一发展颇为出人意料，而后面的故事还会更加精彩。在\hyperref[chap:quantum-cohomology]{第~\ref*{chap:quantum-cohomology} 章}，将研究由量子上同调构造的 \emph{A 模型联络}。特别是，例题~\ref{ex:a-model-yukawa-connection} 将说明：若 Calabi--Yau 三维簇 $V^\circ$ 满足 $h^2(V^\circ)=1$，那么中间上同调 $\bigoplus_pH^{p,p}(V^\circ)$ 有一组基，使 A 模型联络的矩阵\emph{恰好}就是式~\eqref{eq:normalized-yukawa-matrix}，只是 $Y$ 此时为\hyperref[sec:three-point]{第 1.2 节}讨论的 A 模型关联函数。\hyperref[chap:quantum-differential-equations]{第~\ref*{chap:quantum-differential-equations} 章}中的情形更加丰富；例题~\ref{ex:cy-quantum-equation-relation} 从\emph{量子微分方程}的角度讨论 $V^\circ$，并导出算子
\[
\left(\hbar\frac d{dt_1}\right)^2
\left(\frac{\left(\hbar\frac d{dt_1}\right)^2}{Y}\right),
\]
这里 $Y$ 仍是 $V^\circ$ 的 A 模型关联函数。它与 Picard--Fuchs 方程~\eqref{eq:normalized-yukawa-pf} 极为相似（特别是令 $q=e^{2\pi i t_1}$，从而 $\delta=d/dt_1$ 时）。最后，当 Calabi--Yau 三维簇 $V$ 为环面超曲面，$V^\circ$ 为其 Batyrev 镜像时，\hyperref[subsec:toric-mirror-conjecture]{第~\ref*{subsec:toric-mirror-conjecture} 小节}的 Hodge \IndexLocation{indexloc:0795:135}理论\IndexLocation{indexloc:0791:135}环面镜像猜想粗略地断言：在镜像映射给出的坐标变换下，$H^3(V,\mathbb C)$ 上的 Gauss--Manin 联络与 $\bigoplus_pH^{p,p}(V^\circ)$ 上的 A 模型联络一致。一个直接推论就是两侧 Yukawa 耦合的相等。

不过，要达到这一步，还需要做许多工作。首先，必须仔细定义镜像映射。这需要详细研究复结构模及 Kähler 模。复结构模空间有许多紧化，因此可能有许多最大幺幂边界点可供选择；Kähler 一侧的情形同样复杂。\hyperref[chap:moduli]{第~\ref*{chap:moduli} 章}将探究这些问题。

另一个尚缺的要素，是仔细定义上面提到的 A 模型关联函数或三点函数。为此，需要深入研究 Gromov--Witten 不变量。\hyperref[chap:gromov-witten]{第~\ref*{chap:gromov-witten} 章}将讨论这些不变量及其与枚举几何的关系。再将这些内容与\hyperref[chap:quantum-cohomology]{第~\ref*{chap:quantum-cohomology} 章}的量子上同调结合，就具备了陈述\IndexLocation{indexloc:0841:135}镜像对称数学版本所需的全部要素。\hyperref[chap:mirror-theorem]{第~\ref*{chap:mirror-theorem} 章}的主要主题，就是证明其中一些特殊情形。

